% ============================================================
%  H. Høffding, "Kontinuiteten i Kants filosofiske Udviklingsgang"
%  "Continuity in Kant's Philosophical Development"
%  D. Kgl. Danske Vidensk. Selsk. Skr., 6. Række,
%  historisk og filosofisk Afd. IV. 1. Kjøbenhavn 1893.
%  TRANSLATION (English)
%  German, Latin, and French quotations are left in the original,
%  with an English gloss in brackets on first occurrence where helpful.
% ============================================================
\documentclass[12pt,a4paper]{article}
\usepackage[utf8]{inputenc}
\usepackage[T1]{fontenc}
\usepackage{libertinus}
\usepackage{libertinust1math}
\usepackage[english]{babel}
\usepackage[protrusion=true,expansion=true]{microtype}
\usepackage{setspace}
\usepackage[margin=1.2in]{geometry}
\usepackage{fancyhdr}
\usepackage{hyperref}

\hypersetup{
  pdftitle={Continuity in Kant's Philosophical Development},
  pdfauthor={Harald Høffding},
  hidelinks
}

\pagestyle{fancy}
\fancyhf{}
\rhead{Høffding, \emph{Continuity in Kant's Philosophical Development} (1893)}
\cfoot{\thepage}

\setstretch{1.3}

\title{\textbf{Continuity in Kant's Philosophical Development}}
\author{Harald Høffding\\[4pt]
  \small Translated from the Danish}
\date{\emph{D.\ Kgl.\ Danske Vidensk.\ Selsk.\ Skr.}, 6.\ Series,\\
  Historical and Philosophical Section IV.\ 1.\ Copenhagen 1893}

\begin{document}
\maketitle
\thispagestyle{fancy}

\section*{Introduction}

\noindent 1.\ As is well known, Kant's principal work, the «Kritik der
reinen Vernunft» [Critique of Pure Reason] (1781), was preceded by a long
span of years in which the great ideas that were, through that work, to
exert so vast an influence on the intellectual life of the following
century slowly came to maturity. Kant has therefore often been cited as
an example of late maturity. He was fifty-seven years old when the
«Kritik der reinen Vernunft» appeared. But although his significance and
fame are undoubtedly due above all to this work, it would nevertheless be
mistaken to suppose that his preceding writings had value only as
preparatory links in the history of his development and lost their
significance once the principal work had come to light. Just as in their
own time they earned Kant a highly respected name --- so much so that on
their account he was even ranked alongside Lessing, before he stood as the
author of the «Kritik der reinen Vernunft» --- so too do they to this very
day offer more than a merely historical interest. They contain thoughts
whose time is not past. Many of these thoughts Kant --- in part with
metamorphoses conditioned by his later standpoint --- constantly retained
and developed anew in his later works; it is chiefly with these that the
present treatise will be concerned. But in the earlier works there are
also thoughts that in the later works were wrongly cast into shadow or
abandoned altogether. It was not in every respect an advance that Kant
made when he reached his definitive standpoint. With regard to form in
particular, the earlier works have a decided advantage. In them Kant
appears as a clear and elegant writer, which he can by no means be said to
be in all parts of his later works.

It was natural that the new, the grand, and the significant in the
principal work should, both for Kant himself, his disciples, and his
interpreters, cause the earlier works to recede into shadow. Only in
recent years have they been brought forward again out of a purely
historical interest. But this has then been done predominantly in order to
trace the transitional links to the definitive standpoint. This study has
--- together with the Kantian manuscripts that in recent years have for
the first time been brought to light --- brought clarity to Kant's
development in several respects, while at the same time, on certain
points, it has proved impossible to reach any definite result. In this
treatise I shall seek to illuminate some principal points in Kant's
philosophical development, in such a way that I shall particularly
establish the continuity in it and the enduring significance of much in
the earlier works.

Kant himself, on various occasions, declared that his development had had
definite turning points. In a letter to Mendelssohn of 18 August 1783 he
writes of the manner in which the «Kritik der reinen Vernunft» had been
composed: «What was the product of at least twelve years of reflection I
had drafted in the course of four or five months.»\footnote{An entirely
similar statement is already found in Kant's letter to Garve of 7 August
1783: «Den Vortrag der Materien, die ich mehr als 12 Jahre hinter einander
sorgfältig durchgedacht hatte, \ldots\ brachte ich in etwa 4 bis 5
Monathen zu Stande.» (This letter is printed in Alb.\ Stern: \textit{Ueber
die Beziehungen Chr.\ Garve's zu Kant}, p.\ 33 ff.)} This statement points
us back to the year 1769 as the moment of the inception of the train of
thought that led directly to the «Kritik der reinen Vernunft». And with
this agrees what Kant had written far earlier to Lambert (2 September
1770): «A year ago I attained, as I flatter myself, the concept that I do
not fear ever having to alter, though I may have to extend it, and by
which all sorts of metaphysical questions can be tested according to quite
certain and easy criteria, so that it can be decided with certainty how
far they are soluble or not.» In the notes of Kant that Benno Erdmann has
published under the title «Kants Reflexionen», and which contain several
interesting contributions to the history of his development, it is said,
in a fragment that seems to belong to a draft of a projected preface to
the «Kritik der reinen Vernunft»: «Aaret 1769 gav mig meget Lys» [The year
1769 gave me much light].\footnote{\textit{Reflexionen Kants zur kritischen
Philosophie. Aus Kants handschriftlichen Aufzeichnungen herausgegeben von
Benno Erdmann.} II.\ Leipzig 1884.\ p.\ 4.\ (No.\ 4.)} And when one of
Kant's disciples wished to arrange a collection of his minor writings,
Kant did not want works older than 1770 to be included in it (letter to
Tieftrunk of 6 February 1798).

Here, then, we have a date for the beginning of the last period in Kant's
inquiry, and a proof that he himself regarded it as a connected whole. A
turning point in his development is established, and with it the task is
posed --- particularly when one wishes to show the continuity of the
development --- of explaining how this turning point was prepared by, and
hangs together with, the preceding period.

Yet another problem Kant himself has set for those who occupy themselves
with his development, by his reference to the debt he owed to David Hume.
«I freely confess», it is said in the preface to the «Prolegomena», «that
it was the recollection of David Hume which, many years ago, first
interrupted my dogmatic slumber and gave my investigations in the field of
speculative philosophy an entirely different
direction.»\footnote{\textit{Prolegomena zu jeder künftigen Metaphysik.}
Riga 1783.\ p.\ 13.} Here, then, reference is likewise made to an
interruption, in that an impulse from without intervenes and brings about
an entirely new direction in the work of thought. The question then
becomes, first and foremost, at what point in Kant's development this
awakening is to be placed. This question, despite all the learning and
acumen that has been expended upon it, has not reached any unanimous
answer among Kant scholars. There is no point at which there is such a
leap in Kant's development of thought that a definite intervention of an
alien train of thought would have to be regarded as an unavoidable
assumption. Everywhere the progress can very well be explained as a
continuation of the work of the preceding period. One of the most eminent
Kant scholars has supposed that the year 1769, in which, by Kant's express
declaration, the fundamental thought of his definitive philosophy arose,
might also be the moment at which the awakening from dogmatic slumber came
about.\footnote{Friederich Paulsen: \textit{Versuch einer
Entwickelungsgeschichte der Kantischen Erkenntnisstheorie.} Leipzig
1875.\ p.\ 126 ff.} It would indeed undeniably be the simplest thing if
the two turning points mentioned by Kant himself coincided. Whether this
is so can be decided only by a closer investigation of the change that
took place in 1769. In itself there is no likelihood that the two moments
coincide. An awakening is one thing, a discovery another; a new direction
is one thing, a new result another. What happened in the year 1769 may
have been an advance that became possible only because the awakening had
gone before, perhaps long before. The light that dawned for Kant in that
year he may have had to seek for a long time before he found it. And this
is confirmed partly by the fact that in the letters to Garve and
Mendelssohn he speaks of «more than twelve years of work», «at least
twelve years of work», and partly by the fact that in the «Prolegomena»
he states that the distinction he drew [in the «Dissertation» (1770)]
between the elementary concepts of sensibility and of understanding
succeeded only «after long reflection».\footnote{\textit{Prolegomena zu
jeder künftigen Metaphysik.} Riga 1783.\ p.\ 119.}

It would, to be sure, be of no small interest if the moment of the
awakening could be determined with precision. The relation between Kant
and Hume is a principal point in the whole history of modern philosophy.
One might present this history as a great dialogue, in which Hume's
rejoinder and Kant's answer to it stand as the most significant
contributions. Both through the opposition between them and through the
kinship between them, the study of their works is a source of ever-renewed
orientation with respect to questions of philosophical principle. It would
then be of great interest to learn at what point in his development Kant
really and truly heard Hume's rejoinder. That it was not the moment when
he first read Hume lies in his expression: «the recollection of David
Hume». --- It would likewise be of general psychological interest if this
question could be more closely illuminated; we would thereby obtain a
classic example of how a decisive influence can be received in an
eminently independent manner.

But precisely in the difficulties bound up with establishing the moment of
the awakening I see a proof of the continuity in Kant's development.
Precisely because Kant's thinking was in constant unfolding --- precisely
because of his constant striving and tireless self-criticism (at any rate
until the standpoint definitive for him had been reached) --- it is
difficult to decide when and how the external impulses, which on certain
points are known to have been determining, intervened. --- These remarks,
moreover, apply not only to the influence of David Hume, but also to the
influence Kant received from another great contemporary, one which, in
the field of moral philosophy, intervened no less deeply in Kant's
development than Hume's did in the field of epistemology --- namely, that
which came from Jean Jacques Rousseau. Neither would one have become aware
of this had not utterances of Kant in that regard happened to be
preserved.

The difficulty that has shown itself here often recurs in the
investigation of a process of development, especially in the intellectual
domain. To determine the exact relation between inner unfolding and
external influence, where from the very first great forces are in motion,
may easily become an insoluble task. Here it may often seem as if the
inner unfolding could do everything, provided only it received favourable
conditions, without any special external impulse being needed. So it goes
with us precisely in the case of Hume's and Rousseau's influence on Kant;
it is only the external testimonies that lead us to investigate it. Where
Kant scholars disagree with one another, they therefore also employ in the
debate the method of showing how the advance at the particular moment
which the opposing party took to be the moment of the awakening can be
understood without external influence. This method can be carried through
all the more easily because, in the development of a significant mind, an
influence will very often exert itself in an indirect way, in that it will
consist more in the release of one's own hitherto bound power than in a
communication of material. What becomes the result of the awakening need
not then resemble the awakening influence. The word by which one is
awakened often has nothing at all to do with the deed to which one is
awakened.

Although in what follows I occasionally permit myself a conjecture as to
the moment of the awakening, I nevertheless attach the greatest interest
to the dispute about it because it has strengthened the impression of the
continuity in Kant's development. Kant himself, moreover, was likewise not
of the opinion that his pre-critical writings should have entirely lost
their significance. In a draft of a preface to the «Kritik der reinen
Vernunft» he says: «By this treatise of mine the value of my earlier
metaphysical writings is entirely annihilated. Only the correctness of the
idea I shall still try to save.»\footnote{\textit{Reflexionen.} II.\ p.\ 5
(No.\ 7).} A kinship of thought is therefore acknowledged after all,
despite the awakening and the turning point. It is this that is to be
particularly treated in the present treatise.

For the sake of an overview I divide the exposition into four sections, in
such a way that within each of these we investigate the relation between
the earlier and the later standpoint with respect to a single definite
point.

\section*{I.\\ The Concept of Cause}

\noindent 2.\ In Kant's student years his interest was taken up with
Wolff's philosophy and Newton's physics. Toward both he adopted an
independent position, which was a necessity for the very reason that there
was a conflict between them, since Wolff still followed the
Cartesian--Leibnizian physics and contested Newton's doctrine of
attraction. This very conflict was bound to set reflection in motion, and
we accordingly find, already in Kant's first independent period
(1755--62),\footnote{Prior to this period there lie only a treatise on
the estimation of living force (1747) and two short scientific essays
(1754).} attempts by him to alter his philosophical foundation in such a
way that Newton's physics could come into its own. Kant did not stand
alone in the Wolffian school with his recognition of
Newton.\footnote{Kant's teacher, Martin Knutzen, had already sought to
combine Newton's physics with Wolff's philosophy, and he introduced his
pupil to the study of both. Cf.\ Benno Erdmann: \textit{Martin Knutzen
und seine Zeit. Ein Beitrag zur Geschichte der Wolfischen Schule und
insbesondere zur Entwickelungsgeschichte Kants.} Leipzig 1876. --- In
Denmark the Wolffian Jens Kraft, professor at Sorø, declared himself in
his «Cosmologie» (Copenhagen 1752) against the Cartesian physics and for
Newton's.} But he proceeded in a more independent manner than other
Wolffians, and that not only toward Wolff but also toward Newton. In his
famous youthful work «Allgemeine Naturgeschichte und Theorie des Himmels»
[Universal Natural History and Theory of the Heavens] (1755) he takes up
the problem of the natural connection in the universe at the very point
where Newton had not only abandoned it but had even declared it
insoluble. Despite his great reverence for Newton, Kant could not rest
content with his pronouncement that no natural cause could be given for
the fact that the planets move around the sun in concentric orbits, in
the same direction and roughly in the same plane, and that a corresponding
state of affairs obtains with respect to the satellites. Kant regarded an
appeal to a supernatural cause as unscientific, and attempted to show that
the same law-governed order which now unites all bodies and elements into
one great whole had also displayed itself at the origin of the world
systems. This led him to his cosmogonic hypothesis, in which he sought to
show that the origin of the solar system (and, by analogy, of other world
systems) can be understood through the action of forces that are at work
to this very day. It was the principle of actuality, the demand for a
\textit{vera causa} laid down by Newton, which he applied to what is
distant in time, whereas Newton had applied it only to what is distant in
space.

Kant's zeal to carry through a mechanical conception of nature is closely
bound up with a philosophical train of thought that was already developing
in him at that time. It was, he held, a false presupposition that one
proceeded from when one supposed that nature, left to itself, would
produce only disorder and chaos: not by chance, but in accordance with its
own laws, nature produces what is purposive. And precisely the fact that
the various elements thus work together from first to last and are
encompassed by one and the same natural order testifies that in their
essence and origin they are not independent of one another. The mechanical
connection among all things in the world thus leads to the assumption of a
unitary ground for all things. There are two conceptions to which Kant
sets himself in opposition: on the one hand absolute atomism, which
assumes mutually independent atoms that come to work together only by
chance; on the other hand physico-theology, which assumes an external
intervention of the Deity in the world-machinery. The three concepts
element, law, and force he was unwilling to keep separated in an external
manner.\footnote{The intimate belonging-together of the concepts element
and force is treated in particular in a treatise by Kant of the following
year: \textit{Monadologia physica} (1756), in which the atoms (i.e.\ the
absolute atoms) are conceived as points of force, in a manner similar to
that in which, in our own day, Fechner has conceived them in his
«Atomenlehre».} This consideration became of importance especially with
regard to the problem of the origin of the harmony and purposiveness found
in the world. Kant complains that an almost universal prejudice has
prepossessed most investigators against nature's capacity to produce
anything ordered by its own general laws. In his eyes it is precisely the
order and harmony shown in what is produced according to nature's laws
that bear witness that «the essences of all things have a common origin in
one fundamental being». The common connection points to a common
dependence. The more one comes to know nature, the more one is convinced
that the things in the world, in their innermost essence, are not severed
from one another and foreign to one another.

In this youthful work, then, Kant finds the connection between scientific
knowledge and religious faith not by breaking through the connection which
the former finds in nature, but precisely by proceeding from this
connection. And he does not conceive teleology and mechanism as
opposites, but demands that the purposive be explained as the fruit of a
development according to general natural laws.

3.\ The same thought that Kant had applied in practice in his astronomical
hypothesis he expressed at the same time in abstract form in his
habilitation treatise on the principles of
cognition.\footnote{\textit{Principiorum primorum cognitionis metaphysicæ
nova dilucidatio.} (Sämtl.\ Werke, ed.\ Rosenkranz and Schubert.\ I, p.\
1--44.)} He here calls it the principle of coexistence. If one wished ---
such is Kant's train of thought --- to assume that the individual things
(substances) in the world existed each by itself, independent of one
another, one would not be able to infer from any single one of them the
existence of the others. But this is not how things stand. By virtue of
the mutual connection among all things in the world, one can draw
inferences from one to another. This cannot be due to the mere coexistence
of things, but must rest upon a community of origin, upon a common
dependence on one and the same being. In the natural order itself, then,
there lies a proof of a highest cause. This demonstration of a common
source for all things in the world as the presupposition of the
interaction according to general laws Kant believes himself to have been
the first to give (\textit{primus evidentissimis rationibus adstruxisse
mihi videor}).

We stand here before a thought that Kant not only attached great
importance to in his first works, but to which he later returned again and
again, and which is characteristic of his manner of treating problems.
Where Kant stands before a great principle (as here, for example, the
mechanical conception of nature), he takes it in its full range and seeks
either to show that on closer reflection it leads to quite different
consequences than it seemed at first, or else that the whole sphere for
which it holds does not express absolute reality but only one side of it.
The first procedure he employs in his earlier period, the second in his
later period. This procedure has a certain grandeur about it; it is
advantageously distinguished from the so frequent halving of concepts and
from the short-sightedness that would rather seek for itself a brief
shelter in apparent exceptions than resolutely think the principle
through. But it also has its dangers, as Kant's own example will later
show us.

4.\ Seven years after the habilitation treatise and the publication of the
cosmogonic hypothesis, Kant took up the same train of thought that had
animated them, in his work «Einzig möglicher Beweisgrund zu einer
Demonstration des Daseins Gottes» [The Only Possible Argument in Support
of a Demonstration of the Existence of God] (1762). Here the critique of
physico-theology, which had already been initiated in «Allgemeine
Naturgeschichte und Theorie des Himmels», is carried out in more detail.
Kant reproaches it for regarding the harmony and purposiveness in nature
as accidental and therefore seeking a cause for them outside nature. It
thereby becomes unscientific and has often hindered the advance of
knowledge, in that usefulness and perfection were regarded as signs that
the thing which possessed them had no natural origin, so that one
refrained from all further inquiry. And even if one enters into the
physico-theological train of thought, this nevertheless does not lead to
what religious faith demands. For it explains only the ordering of the
elements, not their origin. The Deity whose existence can be inferred
along this line of thought thus becomes merely the architect of the world,
not the world's absolute source.\footnote{It is of interest to note that
Kant's critique of physico-theology dates already from 1755 and 1762. For
in its main features it agrees with Hume's critique in his «Dialogues on
natural religion», which appeared in 1779, two years before the «Kritik
der reinen Vernunft» (although they had been composed many years earlier).
Kant's independence thus stands firm. When he later came to know the
«Dialogues», he set great store by them.}

Nor do the other current proofs of God's existence find favour in Kant's
eyes. Already in the habilitation treatise he had aligned himself with
Crusius' critique of the ontological proof. Crusius had
shown\footnote{Chr.\ Aug.\ Crusius: \textit{Entwurf der nothwendigen
Vernunft-Wahrheiten.} Leipzig 1745.\ §\ 235.\ (Cited here from the 3rd
edition, 1766.)} that this proof, on which the preceding metaphysical
systems (especially those of Descartes and Wolff) had laid so much weight,
rested on a fallacy. Kant carries Crusius' critical train of thought
further in a manner similar to that of the later, well-known development
in the «Kritik der reinen Vernunft» of the difference between concept and
existence. --- Nor does he let the cosmological proof stand. An inference
from the world given in experience cannot lead to the assumption of an
absolutely necessary being.

The critique of natural theology that later, in the «Kritik der reinen
Vernunft», caused such a great sensation, is thus already to be found in
the «Beweisgrund». Were one to confine oneself to Kant's principal works,
one would overlook the significant fact that several of the most important
sections had come into being far earlier.

Yet in 1762 Kant had not given up every attempt at a theoretical grounding
of natural theology. He does not, to be sure, regard proofs as necessary,
since religious faith has another foundation than demonstration. But since
he considers it valuable to test how far one can get along the path of
reason, he sets up the following modification of the ontological proof:
«All possibility presupposes something actual, in and through which
everything thinkable is given. There is, then, an actuality by whose
abolition all inner possibility as a whole would fall away. But that whose
abolition or negation annihilates all possibility is absolutely necessary.
Therefore there is something that exists with absolute
necessity.»\footnote{\textit{Einzig möglicher Beweisgrund zu einer
Demonstration des Daseyns Gottes.} Königsberg 1763 [i.e.\ 1762].\ p.\ 29.
--- This train of thought is already found in the habilitation treatise
(Sectio II.\ Prop.\ 7), only in a more dogmatic form.} This train of
thought is not very clear, and Kant himself offers it not as a
demonstration but only as a ground of proof, the only one he regards as
possible. If it cannot lead to the goal, he says, then none can. --- That
it does not lead to the goal lies in the fact that, even though every
possibility must have its ground in something actual, this actual thing
need not itself be absolutely necessary; it may itself have its ground in
another actual thing --- and so on.

What chiefly interests Kant in this attempted recasting of the ontological
proof is, however, that it would make possible a complete rounding-off of
the train of thought familiar from his earlier writings, which led from
the mechanical connection of nature to a unitary ground for all things.
If each of the things in the world had its own separate necessity, it
would be a matter of chance whether they fitted together in such a way
that a whole could come out of them. For interaction to be possible, the
things must from the very first constitute a system and all be dependent
on a common ground. There must be a principle in which everything
heterogeneous is connected, and in which all the manifold is in unity.
This unitary ground, from which the first possibilities of things spring,
must at the same time be the ground of all wisdom and goodness: only
thereby does it become possible that the law-governed cooperation of
things can accord with what wisdom and goodness demand.

This is the boldest swing that Kant's thinking has taken in a speculative
direction. Far beyond the usual anthropomorphic concept of God, he seeks
back to a foundation from which both the wisdom that proclaims itself in
the world and the elements that cooperate in the world spring. In his zeal
to draw the last consequence of the basic fact from which he proceeds ---
namely, the mechanical connection of nature and the inner connection it
presupposes among all the elements of the world --- his thought finally
loses itself in a mystical abyss. Had he continued to dwell down there
instead of returning to the daylight of rational consciousness, he would,
in the twilight among the few thinkers who moved there, have caught sight
of Jakob Böhme and Spinoza. For when Kant, the first time he introduced
the train of thought whose last consequence he here drew, had with
self-assurance declared himself to be its first originator, this was not
historically correct. In mystical form it is found in Böhme, and in clear
connection with the mechanical conception of nature and as its consequence
it is found in Spinoza, whose entire system is really a carrying-through
of the thought that the causal connection is unintelligible when the
individual things in the world are regarded as self-subsistent and
independent beings (substances in the strictest sense). In his treatise
«de emendatione intellectus» Spinoza sets up precisely the law-governed
connection among phenomena as the basic fact from which thinking is to
take its starting point. Spinoza's one substance corresponds to Kant's
unitary ground. --- Kant never knew Spinoza at first hand, and would
surely have been highly astonished at the agreement into which he had come
with the solitary thinker, for whom the time of understanding had not yet
begun when Kant wrote his habilitation treatise and his «Beweisgrund».

In our own day two ingenious thinkers, in their effort to maintain an
idealist conception of the world without compromising the consequences of
the mechanical conception of nature, have developed a train of thought
similar to that of which Kant believed himself to be the first originator.
In the philosophical views of Fechner and Lotze it plays an important
role. Lotze in particular has, in an acute manner, grounded his idealist
monism on an analysis of the concept of
mechanism.\footnote{\textit{Drei Bücher der Metaphysik.} Leipzig 1879.\
p.\ 135 ff.\ (Already in \textit{Mikrokosmus.} Third Book, ch.\ 5.) In the
most recent times Fr.\ Paulsen has aligned himself with this way of
regarding the matter (\textit{Einleitung in die Philosophie.} Berlin
1892.\ p.\ 212--224).}

5.\ A thought of such depth and interest one could not imagine Kant
letting fall again, once he had conceived it. It accordingly runs, under
various forms, through his later philosophy, although the altered point of
view gives it a different position and form. --- In the same year in which
the «Beweisgrund» came into being, there occurred (as will be shown later)
an important change in Kant's interest. This turned from the objective
problems to an investigation of the method for treating them. It was along
this path that Kant approached the turning point mentioned earlier, in
which the critical philosophy came into being in principle. In the very
work in which this principle was pronounced (\textit{Dissertatio de mundi
sensibilis et intelligibilis forma atque principiis,} 1770) we meet again
(in ch.\ 4) the train of thought familiar from the older writings. Besides
the connection among the things in the world, which appears in the forms
of time and space, there is postulated an objective principle for the real
interaction by which a world-whole becomes possible. The question runs:
«How is it possible that substances can stand in a relation of interaction
to one another and thereby belong to one and the same whole, which we call
world?» (§\ 16). And the answer is that this is possible only when they
all have one source and thus are not substances in the absolute
sense.\footnote{\textit{Totum e substantiis necessariis est impossibile}
(§\ 18). --- \textit{Unitas in conjunctione substantiarum universi est
consectarium dependentiæ omnium ab Uno} (§\ 20).}

When Kant sent the Dissertation to Lambert, he remarked (letter to Lambert
of 2 September 1770) that the chapter in which this train of thought is
found could be passed over as inessential (\textit{unerheblich}). The fact
was that an entirely different problem had now pushed itself into the
foreground, in that Kant had come to the conviction that space and time
are merely subjective forms of intuition, so that things, insofar as we
apprehend them in space and time, are known only as phenomena. He still
supposed, to be sure, that by means of the concept of cause and other
concepts of the understanding we could know things in themselves
(noumena). But immediately after the appearance of the Dissertation he
nevertheless felt great difficulties with this assumption. By pursuing
them he then came (as we shall later show somewhat more closely) to the
result that all our scientific knowledge concerns only phenomena, not
things in themselves.

It is clear that he could now no longer, in the same way as before, infer
from the mechanism of nature to an objective unitary ground. The principle
of the firm connection in the world of phenomena became a purely
subjective principle, the connecting power of consciousness. The point of
unity of the world-view now lay no longer in a mystical primal ground, but
in the knowing subject itself. Thus it is said in a note that seems to
date from the 1770s (the interval between the Dissertation and the «Kritik
der reinen Vernunft»): «The exposition of what is given rests (when one
regards the material as undetermined) on the ground of all relation and of
the linking-together of representations (or sensations). The
linking-together is grounded \ldots\ not on the mere phenomenon, but is a
representation of that inner act of consciousness: to connect
representations, not merely to place them side by side for intuition, but
to make a whole out of them as regards content \ldots\ The exposition of
phenomena is thus the determination of the ground on which the connection
of sensations in them rests.»\footnote{\textit{Lose Blätter aus Kants
Nachlass.} Communicated by Rudolf Reicke.\ First fascicle.\ Königsberg
1889\ p.\ 16.}

With this thought the critical philosophy was completed. It rests on the
fundamental thought that our entire conception of the world bears the
stamp of the mode of operation of our mind. The laws of our reason hold
for all possible phenomena, because these are known to us only in that our
reason intuits and thinks them in its own way and according to its own
laws. But the continuity in Kant's thinking shows itself here, for also
for the critical philosophy the causal connection (besides the connection
in space and time) is the basic fact from which the inference to the
principle of unity is drawn --- only that the principle of unity, as said,
is now subjective, not objective. The fundamental concept of the «Kritik
der reinen Vernunft», synthesis as the expression of cognitive activity at
all its levels, is the subjective counterpart of the objective unitary
ground to which Kant in his youthful works sought back. From seeking the
ground that holds the world together and makes it an objective whole, he
has now passed over to seeking the ground that holds the world-picture
together and makes it a subjective whole. In its innermost essence, as
connecting unitary activity, as synthesis, the knowing mind is then a
prototype of everything that is to be capable of being known, since this
must always have a character of unity. As Kant says: «Ich bin das Original
aller Objecte» [I am the original of all objects], or: «Das Gemüth [i.e.
consciousness] ist sich selbst das Urbild der Synthesis» [The mind is to
itself the archetype of synthesis].\footnote{\textit{Lose Blätter} I.\
p.\ 19.\ 20.} --- No longer substance, but synthesis, is the fundamental
concept. Thereby is designated the critical philosophy's opposition to the
dogmatic.

6.\ In the «Kritik der reinen Vernunft» there is, as regards the critique
of the so-called proofs of God's existence, nothing essentially new in
relation to the pre-critical writings. What is new here is something
negative: the omission of the train of thought which in its time (still in
the Dissertation) had, for Kant, formed a bridge between scientific
knowledge and religious faith. This bridge fell away when the application
of the concepts of the understanding was restricted to the world of
experience. Only by a leap is it now possible to pass from the conditioned
world of experience to the unconditioned that religious faith seeks, and
this leap occurs by virtue of practical, moral motives, not of a craving
for knowledge. Experience, which is always conditioned, cannot ground the
thought of an unconditioned cause.\footnote{\textit{Kritik der reinen
Vernunft} 2nd ed.\ p.\ 665. --- The train of thought by which Kant in his
time (see §\ 4) had attempted to recast the ontological proof he has now
also given up as demonstrative. But the idea of the absolute being as the
ground of all possibilities has not thereby fallen away. It appears as
«the transcendental ideal», the highest measure of everything to which
reality is ascribed. See \textit{Kritik der reinen Vernunft.} 2nd ed.\
p.\ 600; 606. Already in the «Dissertation» the highest theoretical idea
is \textit{exemplar aliquod, omnium aliorum, quoad realitates, mensura
communis} (§\ 9). Here, then, a metamorphosis analogous to the
subjectivising of the «unitary ground» has taken place.} Herewith the
inference from the causal connection in the world to an absolute unitary
ground has fallen away.\footnote{The problem of interaction is, however,
mentioned in \textit{Prolegomena} p.\ 98 f.\ and \textit{Kritik der reinen
Vernunft.} 2nd ed.\ p.\ 265 note, 292 f., 428, and in a manner that shows
that Kant has not forgotten his old thought. Cf.\ also \textit{Reflexionen
Kants.} II.\ p.\ 219 f.\ \textit{Lose Blätter} I.\ p.\ 274.}

But the great interest in causal knowledge and the consistent
carrying-through of it in its strictly scientific form does not fall away
with the restriction of its applicability to the world of phenomena. It is
recovered in the manner in which Kant traces all the various principles
that contain the conditions for a possible experience back to the law of
continuity. The four classes of principles which in Kant's system
correspond to the four classes of categories (just as these in turn
correspond to the four classes of logical judgments) are summed up in four
mottoes: 1) \textit{In Mundo non datur hiatus.} 2) \textit{In Mundo non
datur saltus.} 3) \textit{In Mundo non datur casus.} 4) \textit{In Mundo
non datur fatum.} And as a common characteristic of them all it holds that
they exclude whatever would interrupt the continuous
connection.\footnote{\textit{Kritik der reinen Vernunft,} 2nd ed.\ p.\
282.} Kant has, in this remarkable passage, gathered his entire theory of
knowledge into the law of continuity. Had he laid this point of view at
the foundation instead of following the logical systematics he had to his
delight managed to construct, his fundamental thoughts would have come
more into their own. There are traces that he thought of taking this path,
in that in a note from the 1770s he tries to derive the causal law from
the general law of continuity: «That everything contingent, or whatever
comes into being, must have its ground, comes from the fact that without a
principle there would be no continuity in the phenomena, and without a
rule no identity.»\footnote{\textit{Reflexionen Kants.} I.\ p.\ 308 (No.\
1074). Cf.\ p.\ 514 (Nos.\ 1747--50).}

The thought that here lies at the foundation is that what we seek in the
understanding of a phenomenon is not a mere external juxtaposition of it
with other phenomena, but such a close and definite connection among them
that the phenomenon we seek to understand can stand as a continuation of
the preceding phenomena and together with them form a continuous series.
The incoherent and isolated is unintelligible to us. This important
consideration Kant has only hinted at; but it lies as a germ in the
conception and application of the concept of cause for which he spoke up
from his youth onward. The most acute of his nearest successors has
clearly expressed this thought.\footnote{Salomon Maimón: \textit{Versuch
über die Transscendentalphilosophie.} Berlin 1790.\ p.\ 140: «Die Ursache
der Erscheinung [aufzusuchen, heisst] das Stetige in derselben
aufzusuchen, und die Lücken unserer Wahrnehmungen auszufüllen, um sie
dadurch zu Erfahrungen zu machen. Denn was versteht man sonst in der
Naturlehre unter dem Wort Ursache als die Entwickelung einer Erscheinung
und Auflösung derselben, so dass man zwischen ihr und der vorhergehenden
Erscheinung die gesuchte Stetigkeit finde.»}

The unbroken connection of nature Kant energetically maintains for the
phenomenal world. In the interesting section «Die Disciplin der reinen
Vernunft in Ansehung der Hypothesen» he demands that no grounds of
explanation be employed which do not, according to known laws, hang
together with the given phenomena. It is the principle of \textit{vera
causa}, which he himself had applied in so grand a manner in his cosmogonic
hypothesis. It retains its place in his methodology. And he demands that
it be applied particularly where the greatest temptation to violate it has
shown itself, namely in the explanation of order and purposiveness in
nature. «Nature's order and purposiveness», he says, «must always be
explained by natural grounds and according to natural laws, and here even
the wildest hypotheses, provided only they are natural, are more tolerable
than supernatural ones, that is, than the appeal to a supernatural author,
which one assumes for this need.»\footnote{\textit{Kritik der reinen
Vernunft.} 2nd ed.\ p.\ 800 f.}

7.\ Still more decidedly does Kant, in his last work in the field of
theoretical philosophy, return to the great thought of his youth. The
«Kritik der Urtheilskraft» [Critique of the Power of Judgment] (1790) sets
itself in a way the same task as «Allgemeine Naturgeschichte» and
«Beweisgrund», namely the investigation of how beauty and purposiveness in
the world are to be explained. Kant comes, after many critical
reservations, to the result that the difference we make, in our way of
regarding nature, between the mechanical and the teleological conception
does not justify the assumption that an analogous duality of principles
exists in the essence of things. He seeks to show that when, in the case
of certain natural phenomena, we are led on account of their special
peculiarity to assume a purpose as explanatory ground, while in the case
of simpler phenomena we do not feel this need --- this has its ground in
the nature of our understanding, and we may not ascribe absolute
significance to it. The so-called teleological problem is for Kant only
one form of the general problem that all our experience sets us: how the
manifold can be connected into unity through common laws --- the problem
that has followed Kant through his entire philosophical development. And
although he --- at least as regards mode of expression --- makes as great
concessions as possible to natural theology, he nevertheless points, as
his real result, to the possibility that «the physico-mechanical and the
teleological connection in the ground of nature unknown to us may hang
together in one principle, only that our reason is not in a position to
unite them in such a principle.»\footnote{\textit{Kritik der
Urtheilskraft.} §\ 70 (cf.\ §\ 76, 80, 88).} The thought, familiar from
the youthful writings, of a common foundation for the ultimate
possibilities of things and for the divine wisdom here shows itself by no
means to have been given up, although it comes forward with still greater
caution than before. Kant has described an arc from his first
nature-philosophical essays, through the critical purgatory, to the
profound intimations with which his work of thought ends. Despite the
altered points of view, a continuity in the conception of the world
nevertheless preserves itself, which it is here of peculiar interest to
note.

\section*{II.\\ Analysis and Construction}

\noindent 8.\ We now return to an earlier point in Kant's development in
order to pursue another train of thought in him, just as we have sought to
follow the concept of cause. That the various trains of thought intersect
one another in the course of the development goes without saying; but it
may nevertheless have its advantage to disentangle them each by itself.
The continuity thereby emerges more clearly. And on the most important
points we shall not fail to investigate their interaction.

The year 1762 (and the beginning of the following year) is one of the most
fruitful in Kant's literary life. No fewer than four significant treatises
date from this year, which already by this strong production draws
attention to itself. One of these works we already know: «Einzig möglicher
Beweisgrund». It treated the concept of cause in its significance for
natural science and as a connecting link between the knowledge of nature
and the knowledge of God. It contained, as a particular example, a brief
survey of the cosmogonic hypothesis set up seven years earlier. The other
three writings («Die falsche Spitzfindigkeit der syllogistischen Figuren»
[The False Subtlety of the Four Syllogistic Figures]; «Ueber die
Deutlichkeit der Grundsätze der natürlichen Theologie und Moral» [On the
Distinctness of the Principles of Natural Theology and Morality]; «Versuch
den Begriff der negativen Grössen in die Weltweisheit einzuführen»
[Attempt to Introduce the Concept of Negative Magnitudes into
Philosophy])\footnote{The four writings seem (cf.\ B.\ Erdmann:
\textit{Reflexionen Kants.} I.\ p.\ XVII ff.) to have appeared in the
following order: 1) Spitzfindigkeit; 2) Beweisgrund; 3) Deutlichkeit; 4)
Negative Grössen. --- For the relation among their content of thought the
order of composition and publication is of course not decisive.} discuss
seemingly quite different questions. And yet the investigations carried
out in them became of great importance precisely for the conception of the
very concept on whose significance for natural science and natural
theology Kant had, even in the «Beweisgrund», laid such great weight. In
the year mentioned (whose productivity is perhaps explained precisely
thereby) two trains of thought had moved in Kant's consciousness, whose
collision was bound to confront him with an important problem.

In his critique of the syllogistic figures Kant proceeds from the point of
view that every judgment is a comparison. To judge is to compare a mark
with a thing and, as a result of this comparison, either to ascribe the
thing to it or to deny it of the thing.\footnote{«Etwas als ein Merkmal
mit einem Dinge vergleichen heisst urtheilen» (\textit{Ueber die falsche
Spitzfindigkeit.} §\ 1). --- Cf.\ «Träume eines Geistersehers» [Dreams of
a Spirit-Seer] (1766): «Unsere Vernunftregel gehet nur auf die
Vergleichung nach der Identität und dem Widerspruche». (2.\ Th.\ 3.\
Hauptst.) (Kehrbach's edition, p.\ 64.)} A similar conception of the
nature of thinking lies at the foundation of the treatise «on the
distinctness of the principles». Here the difference between the
philosophical and the mathematical method is impressed upon us.
Philosophy's method is analysis, mathematics' construction. Mathematics
can at once form finished concepts with which one can operate. But in
philosophy, which draws its material from experience, the determination of
the concept becomes finished only when the analysis of the given is
completed. Here, then, is the difficulty: how can one convince oneself
that the analysis has been carried far enough, so that there are no more
marks to discover? In mathematics, therefore, definitions form the
beginning, whereas in philosophy they can come only at the end. The
imperfections of earlier philosophy are derived from the fact that one has
operated with unfinished concepts and embarked on premature constructions.
As an example of a concept with which, in dogmatic philosophy, one has
confidently operated as though it were finished and complete, is mentioned
the concept of spirit (thinking substance), which played so great a role
in Descartes, Leibniz, and Wolff. It was arbitrarily constructed, not
grounded on a carried-through analysis.\footnote{«Träume eines
Geistersehers» is really only a closer development of this example, in
order to show what it leads to, to operate with unfinished concepts. The
concept of spirit (in the spiritualist sense) is here declared a
surreptitious concept. The result of the ingenious work is: «Die
Pneumatologie der Menschen kann ein Lehrbegriff ihrer nothwendigen
Unwissenheit in Absicht auf eine vermuthete Art Wesen genannt werden»
(1.\ Th.\ 4.\ Hauptst.). (Kehrbach's edition, p.\ 43.) Cf.\ herewith
Kant's «Nachricht von der Einrichtung seiner Vorlesungen» (1765--66),
where it is said that empirical psychology does not at all ask whether man
has a soul. --- As one sees, not only the critique of theology, but also
the critique of rational [i.e.\ spiritualist] psychology was essentially
finished at this point of Kant's development.} To carry through the
analysis of complex cognitions Kant declares far more difficult than to
combine simple cognitions by synthesis and build inferences upon them.
Metaphysics is therefore the most difficult of all sciences --- but then
no metaphysics has yet been written either! It will still be a very long
time before one can proceed synthetically in metaphysics. This can happen
only when analysis has helped toward clear and fully developed concepts.

A consequence might here lie near at hand: only such relations are
intelligible in which thinking can, along the path of analysis, lead from
the one term of the relation over to the other. This consequence Kant did
not, however, draw in the treatise «on distinctness». By contrast, the
fourth writing, the treatise on negative magnitudes in philosophy,
acquires its great interest from the fact that this consequence is drawn
with full clarity. In that Kant in this writing impresses the distinction
--- so often overlooked by speculative philosophers of all ages ---
between logical negation and real opposition, he naturally also comes to
investigate such cases where something is annulled because something else
enters in; and here he stands before the causal problem. Along the path of
analysis or comparison one cannot show the necessity of the connection
between the entering of the one and the ceasing of the other, any more than
one can show in general that something happens because something else
happens. Kant lets the problem stand unsolved. Only this much is he
certain of, that the principle of contradiction, to which dogmatism would
trace back all grounding, gives no explanation.\footnote{\textit{Versuch
den Begriff der negativen Grössen in die Weltweisheit einzuführen.} (The
statements on the causal problem are found in the concluding remark.) ---
Cf.\ «Träume» (2.\ Th.\ 3.\ Hauptst.): «Unsere Vernunftregel gehet nur auf
die Vergleichung nach der Identität und dem Widerspruche. So ferne aber
etwas eine Ursache ist, so wird dadurch etwas gesetzt, und es ist also
kein Zusammenhang vermöge der Einstimmung anzutreffen; wie denn auch, wenn
ich eben dasselbe nicht als eine Ursache ansehen will, niemals ein
Widerspruch entspringt». (Kehrbach's edition, p.\ 64.)}

Kant has here posed the causal problem, in part in terms similar to those
in which Hume had already posed it. But Kant had really already had to do
with the causal problem in his earlier writings, most recently in the
«Beweisgrund». What lay at the foundation of the train of thought which
for him connected the knowledge of nature and the knowledge of God, and on
which he believed himself first to have entered (see §\ 3), was the
impossibility of understanding the interaction, the causal connection
among the things in the world, when one did not assume a common ground for
them. So long as they stand in their diversity, the causal relation among
them is unintelligible. What has happened in the remarkable concluding
statements in the treatise on negative magnitudes is therefore really not
the setting up of an entirely new problem, but the transposition of a
problem, hitherto treated in metaphysical-objective form, into
\emph{epistemological-subjective} form. It is understandable that Kant,
now that he had reworked anew the views he had first presented seven years
before, and as he at the same time came upon purely logical and
methodological investigations, was bound, through the collision of these
two trains of thought, to come to see that the causal relation poses not
only an objective but also a subjective problem. Of this year's so intense
and productive work of thought, this transposition is the most important
fruit --- to be sure a fruit that could not yet come to full maturity. The
concept which he himself had applied in such an ingenious way in his
inquiry into nature, and which only recently had built a bridge for him
between the regions of religion and natural science, now suddenly showed
itself to contain so great a problem when one raised the simple question
of how one passes from the one term of the relation it indicates to the
other.

That there was such a connection for Kant between the causal problem in
its metaphysical and in its epistemological form is clearly seen from the
statement in «Träume eines Geistersehers» (in the concluding chapter).
Here attention is once more drawn to the causal problem, and that in
connection with the general problem of the nature of spiritual beings and
their relation to matter. It is shown how the problems begin within
speculation, where one confidently operates with the «fundamental
relations» (cause and effect, substance and action), but that one, when
one continues to philosophise, ends by finding difficulties in the
fundamental relations themselves. And this connection between the
metaphysical and the epistemological problem emerges still more clearly in
the letter to Mendelssohn of 8 April 1766, in that here the general
question of how something can be a cause unfolds out of the more special
one of how a spiritual substance can have the power to act and to be acted
upon in relation to matter.

I find, in accordance with the foregoing, the continuity in Kant's
development at a different point than Friedrich Paulsen. This scholar lays
the weight on the fact that Kant, on closer investigation of the
metaphysical concept of God --- according to which God was to be the
sum-total of all realities --- discovered that there are indeed realities
which stand in a relation of real opposition to one another and thus do
not admit of being united. The treatise on negative magnitudes would then
be a special investigation, called forth by the theological considerations
in the «Beweisgrund».\footnote{\textit{Entwickelungsgeschichte der
Kantischen Erkenntnisstheorie.} p.\ 64 ff.} That there is a connection
between the «Beweisgrund» and the treatise on negative magnitudes I do not
deny. Already in the «Beweisgrund» there is maintained, as Paulsen has
shown, the important difference between logical negation and real
opposition. But the main thing for me is the peculiar and sudden turn that
Kant makes in the concluding remark to the last-named writing. He could
have discussed the problem of negative magnitudes at length and thoroughly
without precisely drawing the consequence he there draws. It is this turn
that I can explain to myself only by the fact that the causal relation was
precisely present to Kant as the most important thing in all knowledge. It
was bound, then, to come to a collision.

9.\ Kant's whole mood at the close of the year so rich in work of thought
was decidedly anti-dogmatic. From confidently operating with hitherto
recognised concepts he had now been led to investigate some of the most
important of these concepts and had found them obscure and unfinished. He
now felt great difficulties with what others --- and he himself hitherto
--- found easy. With irony he turns against «the thorough philosophers,
whose number is constantly increasing», with a request to solve for him
these simple questions at which he had come to a halt. Still more strongly
and more boldly did this mood come forth in «Träume eines Geistersehers»,
which mocks the «builders of air-castles» who construct their edifices out
of surreptitious concepts, and toward whom there is nothing else to do but
to wait until they have finished dreaming.

The strong opposition between construction and analysis itself denotes a
break with the philosophy of the preceding age. Kant demands that an
entirely new foundation be procured before the time for philosophical
systematisation could come. His attention was from now on directed toward
method, toward the investigation of the limits of knowledge. He had surely
already at this time\footnote{As Benno Erdmann has shown in the
introduction to his edition of \textit{Prolegomena,} p.\ LXXVI f.\ and
more fully in \textit{Reflexionen Kants.} II.\ p.\ XXXV--XLVII.} struck out
upon the antinomic method for demonstrating the limit of knowledge, in
that he found such a limit where, concerning the same problem,
self-contradictory propositions could be grounded. «I attempted», he says
later,\footnote{\textit{Reflexionen Kants.} II.\ p.\ 4 (No.\ 4). With this
agrees a statement in the letter to Garve of 21 September 1798, that it was
the cosmological antinomies of pure reason that «first awakened him from
his dogmatic slumber and drove him to the critique of reason in order to
remove the scandal of reason's apparent contradiction with itself». (The
letter is printed in Alb.\ Stern: \textit{Ueber die Beziehungen Garve's zu
Kant.})} «in full earnest to prove propositions and the opposite of them,
not in order to maintain scepticism, but because I suspected that an
illusion of the understanding was present, and in order to discover wherein
it consisted». In a letter to Lambert (of 31 December 1765) he says that
his efforts are chiefly directed toward the peculiar method of
metaphysics. He has now, after several years' deliberation, found a method
which he employs in order to avoid imagined knowledge. In every
investigation he looks to see what he must know in order to be able to
solve the problem, and what degree of knowledge the given permits; his
judgment thereby becomes, to be sure, often more limited, but also more
certain than is usually the case. In agreement with this he states in
«Träume» that metaphysics is the doctrine of the limits of knowledge, and
in «Nachricht von der Einrichtung seiner Vorlesungen» that a «Critik der
Vernunft» is a chief task for him in his treatment of logic. To about the
same time surely belongs also a fragment in which it is said: «Die
Metaphysik ist die Kritik der menschlichen Vernunft \ldots\ [sie] ist
subjectiv und problematisch».\footnote{\textit{Reflexionen Kants.} II.\
p.\ 158 (No.\ 507).}

10.\ Those Kant scholars are surely right who have sought to show that the
development that takes place in Kant's thinking in 1762 and the following
years does not necessarily presuppose an influence from without, but in
itself can very well be understood from his preceding development. But on
the other hand I know no period in Kant's development to which the
expression «awakening from dogmatic slumber» would fit so well as here.
Even Benno Erdmann, who places the awakening far later (so much later that
it seems to me he comes into conflict with Kant's statement that the
awakening took place «vor vielen Jahren» [many years ago]), declares the
time in the 1760s to be the period in which «the Kantian thoughts were
most in flux». To be awakened from dogmatic slumber means precisely that
flux comes into the hitherto fixed thoughts. To be sure, if by the
expression «awakening from dogmatic slumber» one will understand the
complete transition to the critical philosophy, then it does not fit this
moment. In this strict sense B.\ Erdmann takes it, when he first finds it
applicable where the hope of an intellectual knowledge of things in
themselves «was not merely eradicated to the last fibre, but was even
replaced by a contrarily opposed
conception».\footnote{\textit{Einleitung zu Prolegomena,} p.\ XCII.} This
is surely to demand too much of an awakening. Kant himself has, in a draft
of a preface to the «Kritik der reinen Vernunft», stated: «It was long
before I found the \emph{whole} dogmatic theory dialectical. But I sought
for something certain, if not with respect to the object, then at least
with respect to the nature and limits of knowledge.»\footnote{\textit{Reflexionen
Kants.} II.\ p.\ 4 (No.\ 3). --- I lay the emphasis on the word «whole».}
He who seeks thus is awake. We have also heard that he himself at this time
regarded the dogmatists as dreaming. I cannot see otherwise than that the
sharp opposition between construction and analysis, the great weight on
method and on the limits of knowledge, as well as the low opinion of what
had hitherto been achieved in philosophy, decidedly place Kant outside the
circle of the dogmatists already at this time. Only from a (in the Kantian
sense) «hypercritical» standpoint could he in this period be called resting
in dogmatic slumber. Should Kant nevertheless in the year 1783 have
reckoned the period of the 1760s to dogmatism, then he has done himself a
great injustice. His own definition of the concept of dogmatism was, after
all, that it consisted in «the presumption of getting through merely with
a pure knowledge of concepts, according to such principles as reason has
long had in use, without investigating the manner in which, and the right
by which, it has arrived at them»,\footnote{\textit{Kritik der reinen
Vernunft.} 2nd ed.\ Preface, p.\ XXXV.} or, as it is said in another
place: «the general confidence in the principles of metaphysics without a
preceding critique of the faculty of reason itself».\footnote{\textit{Ueber
eine Entdeckung, nach der alle neue Critik der reinen Vernunft durch eine
ältere entbehrlich gemacht werden soll.} Königsberg 1790.\ p.\ 78.}
According to this definition Kant could not, in the period from 1762, be
called a dogmatist. He had, to be sure, not --- not even in «Träume eines
Geistersehers» --- given up the hope that diligent work along the path of
analysis might one day make a constructive knowledge possible. But he had
a clear eye for the difficulties and very little confidence in the current
principles, as well as a definite conviction that the result, if any came,
would be bound up with a limitation of our knowledge. If he himself later
reckoned this period to the period of dogmatic slumber, then the results of
the «Critique» must have blinded his view of his own past. --- The question
of the moment of the awakening it is nevertheless not possible to decide
with certainty according to what is available. I can only not see otherwise
than that the placing of it in the year 1762--63\footnote{To this moment
the «awakening» is assigned by Kuno Fischer, Riehl, and Vaihinger (the
last of whom, however, also assumes a later awakening). My grounds are,
however, different from those of these scholars.} is most probable, though
such eminent Kant scholars as Fr.\ Paulsen and Benno Erdmann have come to
other results. As repeatedly remarked, the uncertainty about this question
is a testimony to the continuity in Kant's development.

11.\ Benno Erdmann\footnote{«Kant und Hume um 1762». \textit{Archiv für
die Geschichte der Philosophie.} I.\ p.\ 62--77; 216--230.} believes, from
Herder's utterances from the years in which he was Kant's auditor
(1762--64), that one can infer that Kant cannot at that time have been
awakened by Hume. One would otherwise, Erdmann holds, have noticed
something, in the manner in which Herder speaks of Hume, of the enthusiasm
which Kant must undoubtedly have instilled in his auditors for the English
thinker.

This assumption does not seem to me compelling, especially when one holds
fast that the awakening must have had a highly indirect character, so that
it essentially consisted in a release and full clarification of thoughts
and doubts that had already begun to work their way forward. It was, by
Kant's statement, «the recollection of Hume» that awakened him. So Kant
had earlier read Hume (it is the latter's «Inquiry» in German translation
that is in question here), but the thought of Hume's posing of the problem
influenced him only later. It is easily understandable that this
recollection of a reading which in its time had made no stronger
impression could come forth in Kant precisely in the year when the two
trains of thought that had hitherto run side by side in his consciousness
--- the causal relation as expression of the connection of various
elements, and thinking as analysis, operating with the principle of
contradiction alone --- were both taken up for new, penetrating treatment
and easily had to come into collision. One may suppose that «the
recollection of Hume» either favoured this collision or strengthened its
effects. But Kant can in any case scarcely have been able right away to
have full consciousness of the significance of what was going on in him; a
such consciousness could come only when the process of development that
was initiated by the awakening lay quite finished before the retrospective
gaze. Thereby it is also understood that he does not mention the awakening
by Hume in his writings until long afterward and, as Erdmann himself has
shown, on a special occasion.

There is no reason to assume that Kant, in his lectures at the time Herder
heard him, should have spoken of Hume (as theoretical philosopher)
otherwise than he constantly spoke of him later, namely as the sceptic who
forms a sound counterweight and a beneficial ferment over against
dogmatism. Kant moreover still lacked the «higher unity», criticism, which
he later used to let triumph over dogmatism and scepticism, and he had
then to regard Hume predominantly as an opponent to be overcome. In the
utterances that Erdmann cites from Herder, the latter speaks of Hume
precisely in the manner we should expect, namely as a fine mind, but at
the same time as an arch-doubter; he even accuses him of «ein erstrebtes
Zweifeln» [a contrived doubting]. I do not understand how Erdmann, after
these utterances, can infer\footnote{\textit{Archiv.} I.\ p.\ 216.} that
Kant --- when one relies on Herder's testimony --- in these years did not
appreciate in Hume the metaphysical doubter, but only the moralist. The
doubter in Hume Herder does after all seem to have been made sufficiently
aware of! Herder has surely even exaggerated the matter. The expression
«erstrebt» must probably stand to his account, not to Kant's.\footnote{A
letter from Kant to Herder of 9 May 1767 (of which I became aware only
through a kind gift from the librarian Dr.\ R.\ Reicke in Königsberg)
shows that Kant nevertheless at this time had a more positive conception
of Hume than one would suppose from Herder's utterances. Kant wishes his
young friend that he may attain that (by no means unfeeling) calm in its
stead, which is peculiar to the philosopher as opposed to the mystic, and
adds thereto: «Ich hoffe diese Epoche Ihres Genies aus demjenigen was ich
von Ihnen kenne mit Zuversicht, eine Gemüthsverfassung, die dem so sie
besitzt und der Welt unter allen am nützlichsten ist, worin Montaigne den
untersten und Hume so viel ich weiss den obersten Platz einnehmen». This,
in Kant's mouth, great praise must, according to the whole context, apply
not merely to «the moralist». (The letter is printed in
\textit{Altpreuss.\ Monatsschrift.} Bd.\ 28.\ Heft 3 u.\ 4.\ 1891:
\textit{Zu Herders Briefwechsel. Von Victor Diederichs.})} It was, given
Herder's whole natural disposition and intellectual direction, not strange
that Hume's doubt could seem to him exaggerated and arbitrary. Herder had
not the use for it that Kant had, who got his thoughts more strongly set
in motion by it; indeed, Herder could perhaps not even understand what use
Kant could have for it; his later relation to Kant gives ground to suppose
so. What Herder appreciated to so high a degree in Kant at this time was
(according to the well-known statement in the «Letters on Humanity») the
open sense for everything human, as well as for nature, the lively mind
and the clear understanding. For the seeking thinker, wandering along
solitary and thorny paths, Herder had no sense.

12.\ By the clear distinction between construction and analysis Kant had
made a great advance. The fault of the dogmatic systems was that they
proceeded too quickly to construction, even though they naturally could
not entirely dispense with analysis, which --- however incomplete and
unfinished --- in reality always forms the preparation. Kant saw for the
present, in the more exact analysis, the most important means of progress.
And --- what is of particular interest for holding fast the continuity in
his development --- he really never gave up, not even in the «Kritik der
reinen Vernunft», the relation between analysis and construction that he
had presented in the writing «on the distinctness of the principles».

By his emphasis on the significance of analysis Kant meets with Lambert's
efforts. But while Kant sees construction lying in the indeterminate
distance and holds that for the present it is the time of analysis, Lambert
assumes that the goal lies nearer. By comparison and combination of the
simple concepts to which analysis has led, he will without further ado form
axioms and postulates (letter to Kant, November 1765). He sees well that a
closer testing of the various combinations is necessary;\footnote{J.\ H.\
Lambert: \textit{Neues Organon.} Leipzig 1764.\ I, p.\ 323 f.; 499.} but he
feels no need to find special principles according to which this testing of
the possible principles could take place. He approaches it in the
remarkable statement (letter to Kant, February 1766) in which he raises the
question whether the knowledge of the form of our knowledge leads to
knowledge of the content of our knowledge. But the fundamental difficulty
in the transition from analysis of the fundamental concepts to construction
of the principles he has not come to a halt at, nor did he carry analysis
in the definite direction by which Kant was led to find the sole
possibility of constructive knowledge in philosophy. Although Kant
therefore honoured Lambert as his predecessor so highly that he would have
dedicated the «Kritik der reinen Vernunft» to him had he not died before
the work's completion,\footnote{See the draft of a dedication:
\textit{Reflexionen Kants.} II.\ p.\ 1 f.} it was nevertheless his
judgment of him that he moved within mere analysis, without reaching the
«critical» standpoint.\footnote{\textit{Reflexionen Kants.} II.\ p.\ 67 f.\
(Nos.\ 225, 227, 231).}

The analysis that plays so great a role in Kant's definitive theory of
knowledge concerns the manner in which the faculty of knowledge itself
operates. This direction Kant's analysis already took when a more careful
determination of the limits of knowledge proved necessary (see §\ 9). But
it led to decisive results only when, from 1769 onward, he began to
distinguish between the material and the form of knowledge. How this
turning point came about will be discussed later; here it shall only be
investigated of what kind the analysis was by which Kant grounded this
distinction so important for his definitive system.

Kant has nowhere methodically carried out this analysis. He constantly
presupposes it or hints at it in great brevity. In the «Dissertation» it
is said that by attending to the manner in which consciousness operates on
the occasion of experience (\textit{attendendo ad ejus actiones occasione
experientiæ}) we discover the laws according to which this operation goes
on. Now when, in sensible knowledge, a distinction is made between material
and form, such that form means the law for consciousness's ordering of the
received material, then it emerges from the expressions Kant uses that form
is the constant, unchangeable element in sensation, whereas the material is
that which varies infinitely. Analysis thus finds a difference between a
constant and a changing, a determinate and an
indeterminate.\footnote{\textit{Tempus est \ldots\ subjectiva conditio per
naturam mentis humanæ necessaria, \ldots\ quælibet sensibilia certa lege
sibi coordinandi} (Diss.\ §\ 14, 5). --- \textit{Spatium est \ldots\
subjectivum et ideale e natura mentis stabili lege proficiscens, veluti
schema, omnia omnino externe sensa sibi coordinandi.} (Diss.\ §\ 15 D).}
The form stands as an unchangeable prototype (\textit{typus immutabilis}),
which is constantly imitated anew, every time a new content is to be
ordered. In notes from the 1770s the same trait is found. «The exposition
of what is thought rests merely on consciousness; but the exposition of
what is given rests --- when one regards the material as indeterminate ---
on the ground of all relation and connection among the representations».
And when we make clear to ourselves according to what laws this connection
goes on, we obtain «an original of all objects», «an archetype for every
possible synthesis», which can be known apart from, or «independently of
the particular composition in which it was given» (\textit{unabhängig von
der Einzelnheit darin es gegeben war}).\footnote{\textit{Lose Blätter.} I.\
p.\ 16; 19--21. (Cf.\ above §\ 5, conclusion).}

Also in the «Kritik der reinen Vernunft» there are traces of the analysis
by which form is separated from material. When Kant --- by an unfortunate
mode of expression --- says that we can know the form prior to actual
observation, he presupposes that such an analysis has already been carried
out, and by actual observation he means: any definite, particular
observation. --- A few examples will show this.

«When I separate off from the representation of a body what the
understanding thinks about it (substance, force, divisibility, etc.), and
likewise that part of it which belongs to sensation (impenetrability,
hardness, colour, etc.), then I still have something left over from this
empirical intuition, namely extension and figure. These belong to the pure
intuition that takes place in consciousness even without any actual object
of sensation or feeling, as the mere form of sensibility». There is
described here an analysis, that is, a successive directing of attention to
the various sides or elements of the given, and the results of the analysis
are divided into three different classes, which are referred to the three
faculties, understanding, form of intuition, and sensation.\footnote{Kant
expresses himself here in the manner of the old theory of abstraction,
although in other places he definitely rejects this. See «Kritik der reinen
Vernunft». 2nd ed.\ p.\ 457 (the note to the first Antithesis).
\textit{Reflexionen Kants.} II.\ p.\ 126 (No.\ 412).} How the
«separating-off» is more closely grounded is seen from other places.
«Space and time are [sensibility's] pure forms, sensation in general is the
material. \ldots\ The former hang together entirely necessarily with our
sensibility, of whatever kind our sensations may be; the latter may be very
different». «We call the synthesis of the manifold in the imagination
transcendental when, apart from any difference in what is intuited (\textit{ohne
Unterschied der Anschauungen}), it bears on nothing but the connection of
the manifold a priori». «Apperception, and with it thinking, precedes all
possible definite arrangement of the representations.»\footnote{\textit{Kritik
der reinen Vernunft.} 2nd ed.\ p.\ 35; 60; 118; 345.} --- And in a place in
«Prolegomena», where Kant reports on the genesis of the doctrine of the
categories, he says: «I looked about for an act of understanding that
contained every other and offered differences only through the various
modifications or moments under which the manifold of representation is
brought under the unity of thinking. Then I found that this act of
understanding consisted in judging.»\footnote{\textit{Prolegomena.} Riga
1783.\ p.\ 109.}

There arose at the same time the task for Kant of finding common marks for
the constant elements which he could thus gradually single out in our
knowledge. This part of the analysis went hand in hand with the
demonstration of the constant elements, in that Kant first became properly
aware of these when he had found a common fundamental determination for all
our knowledge. This fundamental determination was connection, synthesis,
and since it is a function, Kant regarded it as self-evident that it could
not be given with the mere impressions. But of this side of the analysis we
shall treat more closely in ch.\ 4. Here we confine ourselves merely to the
general significance that the analytic method continued to have for Kant.

This analysis, which aims at finding the very manner in which the faculty of
knowledge operates, is expressly distinguished from «the usual procedure in
philosophical investigations: to analyse the concepts that present
themselves according to their content and bring them to distinctness». The
real task of the critical philosophy is an analysis of the faculty of
knowledge itself, to find the fundamental concepts there where they spring
up, and to demonstrate their application. And this critical analysis, which
separates form from material, cannot proceed «until long practice has made
us attentive and has put us in a position to single out what our own
faculty of knowledge contributes from what we receive through sense
impression». Therefore Kant has nothing against granting to the
mathematician Kästner that the concept of space is abstracted from sensible
representations: «for without application of our sensible mode of
representation to actual objects of the senses, even that which a priori is
contained in these [i.e.\ in the senses] would not at all become known to
us».\footnote{Dissert.\ §\ 15 Coroll. --- \textit{Kritik der reinen
Vernunft.} 2nd ed.\ p.\ 1 f.; 90. --- Wilh.\ Dilthey: \textit{Aus den
Rostocker Kanthandschriften. Archiv für die Geschichte der Philosophie.}
III.\ p.\ 87.}

13.\ The foundation found by analysis for the critical doctrine of
knowledge emerges more strongly in the first edition of the «Kritik der
reinen Vernunft» than in the Prolegomena and the second
edition.\footnote{This has been shown by Benno Erdmann: \textit{Einleitung
zu Prolegomena,} p.\ XXXII--XXXVIII.} There he develops, in a
psychologically very interesting manner, the concept of synthesis as the
general form of conscious activity at the various levels. Later he contents
himself with referring quite generally to the connecting process,
particularly in the domain of the understanding.\footnote{Cf.\ also the
letter to Tieftrunk of 11 December 1797 on the concept of composition as
that which lies at the foundation of all categories.} On account of the
misunderstandings in a psychological-idealist direction of which the
«Kritik der reinen Vernunft» at once became the object, Kant in the later
reworkings sought to divert attention from the psychological analysis of
the cognitive process (from what he in the preface to the first edition
called the subjective deduction), which had only a preparatory
significance, in order that his chief task --- the construction of the
rational propositions that hold for all phenomena and yet are independent
of them --- might emerge all the more sharply.

But the sharp opposition between the method of mathematics and that of
philosophy, which Kant had presented in the treatise on the distinctness of
the principles, he nevertheless constantly holds fast. In the interesting
section «Die Disciplin der reinen Vernunft im dogmatischen Gebrauche» it is
expressed in part with the same turns of phrase as in the nearly twenty
years older writing. Here too it is shown that the philosophical concepts
gained by analysis, not, like the mathematical, by direct construction,
cannot in the strictest sense be defined, since there can be no apodictic
certainty that the content of the concept is complete. As examples are
mentioned the concepts cause, right, and equity. There may be obscure
accessory representations present which we do not become aware of in the
analysis, and which nevertheless, without our noticing it, play a part in
the application.\footnote{\textit{Kritik der reinen Vernunft.} 2nd ed.\
p.\ 756. --- Cf.\ herewith «Träume eines Geistersehers». 1.\ Th.\ 1.\
Hauptst.\ (Kehrbach's edition p.\ 6 note), where the description of
surreptitious concepts quite recalls the cited passage of the «Critique».}

The difference between the older writings (1762 ff.) and the «Kritik der
reinen Vernunft» is that whereas Kant earlier distinguished between two
kinds of concepts, empirical and mathematical, of which the former are
gained by analysis, the latter by direct construction, he later
distinguishes three kinds of concepts: 1) empirical, 2)
a-priori-discursive, 3) a-priori-intuitive. The first two kinds are gained
by analysis, the last by direct construction. But the analysis by which the
first and the second kind are gained is, as has been shown in the foregoing
(see above §\ 12), of different character.

14.\ For the critique of Kant's theory of knowledge this connection between
the older and the later standpoint is of great interest. For it suggests
the question whether Kant has really discovered an entirely new class of
concepts, which would stand between the empirical and the constructed. In
deciding this question one need not go beyond Kant's own statements. For he
declares in plain words that the a-priori-discursive concepts (i.e.\ the
categories) are not fully finished --- and yet he operates with them to set
up principles, as though they were fully finished! If we look to their mode
of genesis, then they are subject to the same conditions as the ordinary
empirical concepts. They are due, as we have seen, to the observation of
constant elements amid the variation of our experiences. The assumption
that we here have before us the very nature of the faculty of knowledge
(«the sources», as Kant is wont to say) is an inference from, or an
interpretation of, this constancy. It is thus a hypothesis. In the preface
to the first edition of the Critique of Reason Kant also says that «the
subjective deduction» (i.e.\ the demonstration of the sources of knowledge
by means of analysis) might look like a hypothesis, since it aims at «the
demonstration of a cause for a given effect»; but he holds that this is not
so, and promises on another occasion to establish it. This promise he has,
however, not kept, whereas in his later presentations he has considerably
reduced this analytic part.

Kant is thus, in a certain sense, on his definitive standpoint, no further
advanced than he was in 1762. He has not shown the possibility of
determining the «forms» in such a way that there is full guarantee of the
exactness and completeness of the determination. He cannot then really know
himself secure against the same charge that he directed against Leibniz and
Wolff: of operating with unfinished and surreptitious concepts. And there
is a considerable hypothetical element in the system of forms which he sets
up as the foundation for our exact experience. He is right that if we could
be quite certain that these forms must play a part in all our experience,
all our apprehension, then we would also have an a priori knowledge of the
objects of experience. But precisely this full certainty cannot be
established.

For the theory of knowledge there is here no other path to go than to
operate with the foundation found by analysis as with a hypothesis, leaving
it to continued analysis to correct the foundation from which one proceeds.
The fact is that the various operations stand in a reciprocal relation to
one another. One cannot lay aside all possible construction and set about
analysing; that would be to the detriment of analysis itself, since
construction may be necessary as an auxiliary process and experiment.
Kant's own antinomic method was really attempts at constructions that were
to illuminate the character of the presuppositions when these could not be
directly analysed. --- The reverse also holds, of course, that construction
presupposes analysis, for which reason even the mathematical concepts that
seem set up by direct construction prove, on closer investigation, to
presuppose a preceding analysis. Kant's opposition between mathematics and
philosophy is thus not tenable in the form in which he has set it up.

The certainty that in the constant elements of experience we have the
necessary forms of our conscious activity can, however far analysis be
carried, become only approximate, and the inferences that can be built on
this certainty must of course share its fate.\footnote{Cf.\ already
\textit{Ænesidemus} [by G.\ E.\ Schulze]. 1792 p.\ 401; 406.} Kant's
tendency, on his definitive standpoint, to push analysis aside has made him
a dogmatist against his will. He saw more correctly in 1762, when he used
the analytic method as a weapon against the dogmatists.

15.\ That analysis has been a stepchild in Kant's definitive philosophy,
although it has its definite place in it, shows itself especially in the
fact that the derivation of the «forms» from the concrete connection in
which, according to Kant, they first display themselves, is not carried
through. The relation between the two propositions to which Kant equally
subscribes --- «All knowledge begins with experience», and «Not all
knowledge springs from experience» --- he has not clearly demonstrated.
This would only have happened had he presented the various stages of
development through which the «forms» pass from the innate foundation to the
conscious clarity acquired by practice.\footnote{Cf.\ the distinction
between innateness and original acquisition in Dissert.\ §\ 8; Coroll. ---
\textit{Ueber eine Entdeckung} etc.\ p.\ 68. --- \textit{Kritik der reinen
Vernunft.} 2nd ed.\ §\ 27. --- Letter to Hertz of 21 February 1772. ---
From a purely psychological standpoint, too, objection must be raised
against Kant's distinction between material and form. Sensations are not
merely passively received material, but every sensation is determined in
its arising and character by the whole state of consciousness. Here, as far
as the sensations themselves are concerned, a forming and connecting, a
synthesis, takes place. Cf.\ my \textit{Psychology}, 3rd ed.\ p.\ 130. This
objection does not abolish Kant's distinction, but carries it further and
in another manner than he did.}

But not only does Kant presuppose the forms as fully finished; he also
presupposes them in an ideal perfection and purity in which they do not
occur in any actual intuition or experience. That this is so as regards
space and time needs no closer demonstration. Newton's absolute space and
time are with Kant, so to speak, struck inward, have become subjective
forms --- without, however, losing their absoluteness. Thus each of the
forms of intuition is designated in the Dissertation as \textit{typus
immutabilis} (§\ 16 Coroll.). Kant thus imagines prototypes, ideal frames,
to which everything that appears to us must be referred in order to be
worked into a general connection.\footnote{As regards time, Kant has not
distinguished between sensation of time, representation of time, and
intuition of time. Cf.\ my \textit{Psychology.} 3rd ed.\ p.\ 206--214 and
my treatise: «Lotze og den svenske Filosofi». \textit{Nordisk Tidsskrift,}
publ.\ by the Letterstedt Association 1890.\ p.\ 6--10. --- As regards
space, Kant expresses himself in other writings far more clearly than in
the «Kritik der reinen Vernunft». Thus it is said in «Metaphysische
Anfangsgründe der Naturwissenschaft» (Riga 1786).\ p.\ 3 f.: «Der absolute
Raum ist an sich nichts und gar kein Object, sondern bedeutet nur einen
jeden andern relativen Raum, den ich mir äusser dem gegebenen jederzeit
denken kann». --- p.\ 16: «Der absolute Raum ist für alle mögliche
Erfahrung nichts».} A similar prototype (\textit{exemplar}) is, according
to the Dissertation (§\ 9), formed also for intellectual knowledge and
provides the common measure for all reality. In the «Kritik der reinen
Vernunft» the very concept of experience or of nature arises by a similar
ideal construction. The parallel with that «unchangeable prototype» emerges
clearly when it is said: «There is only one experience, in which all
observations are represented in thoroughgoing and law-governed connection,
just as there is only one space and one time; in the former all phenomena
and all relations of being and non-being find their
place.»\footnote{\textit{Kritik der reinen Vernunft.} 1st ed.\ p.\ 110.}
«All our cognitions lie within all possible totality». «All phenomena lie
in one nature and must lie therein, because without the a priori unity
there would be no unity of experience and thus no determination of the
objects in it possible.»\footnote{\textit{Kritik der reinen Vernunft.} 2nd
ed.\ p.\ 185; 263.} --- The last ground that Kant gives for this unity of
experience (or of nature) is that conscious knowledge is possible only when
the given elements are united by synthesis. Synthetic unity is the
principle of all understanding. «If every single representation were quite
foreign to, isolated and severed from the others, then there could never
arise anything such as knowledge is. For knowledge is a whole of compared
and connected representations.»\footnote{\textit{Kritik der reinen
Vernunft.} 1st ed.\ p.\ 97 (cf.\ 2nd ed.\ p.\ 103).} We have therefore, as
we have already shown (§\ 5 concl.; §\ 12), in the concept of the synthetic
character of our knowing consciousness an archetype for everything we can
know.

But to what degree do the actual observations now satisfy the ideal demand
that is, on the basis of this prototype, made of them? About this, plainly,
nothing can be known beforehand. A prototype formed by abstraction and
idealisation of our cognitive activity shows its significance by setting
thought in motion. But whether thought finds adequate counterparts of the
prototype, and whether it can carry the prototype through amid the
unsurveyable manifold of observation, that is the great question, which is
not answered by the demonstration that we have necessary knowledge only if
this happens.

Sometimes Kant also expresses himself, then, as though by the prototype
merely a search were initiated. In the Dissertation, accordingly, the
principle that everything that happens, happens according to the order of
nature, is presented as a subjective principle for our inquiry
(\textit{principium convenientiæ}), on a par with the law of parsimony (§\
30). In a note from the 1770s there is talk of a presumption, according to
which everything can be determined by a rule. In a note «from the first
time of criticism» it is asked: «Since we do not see the possibility of a
real ground, how can we then a priori see that there absolutely must be
such? \ldots\ Does this principle not hold as anticipation, because without
a rule we would have no experiences, and this rule consists in the ordering
of time and space according to general laws?» In a note from the 1780s it
is said: «The principles for the exposition of phenomena are principles for
their intellection, not for their perspicience, for seeking causes, not for
determining them». And in the «Kritik der reinen Vernunft»: «The
understanding is constantly occupied with scrutinising the phenomena with
the aim of finding out a rule for them». «The understanding can a priori
yield no more than to anticipate the form of a possible experience in
general.»\footnote{\textit{Lose Blätter.} I.\ p.\ 35; 136. ---
\textit{Reflexionen Kants.} II.\ p.\ 307 (No.\ 1071). --- \textit{Kritik
der reinen Vernunft} 1st ed.\ p.\ 126; 2nd ed.\ p.\ 303.}

Had Kant remained here, he would not have asserted more than he could
prove. But then he would, to be sure, have had no ground to set up his
sharp difference between category and idea, since it would then show itself
that the category already poses ideal demands which real knowledge always
satisfies only imperfectly.\footnote{Already Salomon Maimón demonstrated
this. «Für mich ist Erfahrung im strengen Sinne kein in der Anschauung
darstellbarer Begriff, sondern eine Idee». \textit{Kritische
Untersuchungen.} Leipzig 1797.\ p.\ 154. --- «Die kritische Philosophie
kann nur hypothetisch philosophiren». \textit{Die Kategorien des
Aristoteles.} Berlin 1794.\ p.\ 133.} Kant was carried beyond the goal by
his confidence in the «Copernican» thought, in which (as will be mentioned
later) the discovery of 1769 consisted. It gave him the courage to
undertake the synthetic construction which only a few years before had
stood for him as a distant goal. From the archetype and the presumption he
passed without further ado to the law, as though this could be derived from
those. He regards real knowledge as given with the concept of knowledge:
the old ontological train of thought has here outwitted its great opponent!
The transcendental deduction, which had cost him so much trouble, and of
which he was so proud, was in reality epistemological alchemy. But just as
the alchemists often, in their search for the philosopher's stone,
discovered along the way new chemical relations, so Kant, in his eager
striving to give an absolute solution of Hume's problem, came to make
valuable contributions to the understanding of the nature of our knowledge.

16.\ Only a single point shall still be brought forward to illuminate
Kant's analytic method. He proceeds, as we have seen (§\ 12), from the
assumption that the constant elements in our experience must derive merely
from our cognitive activity itself and be only an expression of its laws.
Had Kant gone more fully into the analysis that led him to set up the
forms, he would also have had to come to test more closely the
justification of this assumption. He proceeds, after all, from the very
first from the assumption that consciousness, the knowing subject, is not
the only existing thing, but is only a member of the whole of existence and
stands in a reciprocal relation with the rest of what exists. The manner in
which the subject apprehends existence can then not be determined merely by
its own nature, but must also be determined by the nature of what exists.

There are several unnoticed intimations in Kant in this direction. In the
«Kritik der reinen Vernunft», first edition (p.\ 358), there is talk of
«the something that lies at the foundation of the outer phenomena and
affects our sense in such a way that it obtains the representations of
space, matter, figure, etc.». And on the occasion of the question how
intuition of space is at all possible in a thinking subject, it is remarked
that to this only this much can be said, that the representations of outer
phenomena have their cause in the thing in itself («the transcendental
object») (ibid.\ p.\ 393). The manner in which I apprehend a thing need
not, after all, at all resemble the very thing that calls forth the
apprehension (the sensation or the spatial intuition) in me --- it is said
in Prolegomena p.\ 64. Spatial intuition is, it is said in the second
edition of the Critique of Reason (p.\ 66, 69), determined by the relation
of the object to the subject. The clearest statements are found in some
places in the «Dialectic». It is impossible, it is said, to answer the
question «why the transcendental object gives our sensible intuition
precisely only spatial intuition and not another kind of intuition». «Many
forces in nature, which manifest their existence by certain effects, are
and remain unfathomable for us; for we cannot trace them far enough by
observation. The transcendental object lying at the foundation of the
phenomena, and with it the ground of why our sensibility has precisely
these definite highest conditions and not others, are and remain
unfathomable for us, although the matter itself is otherwise given, only
not seen through» (2nd ed.\ p.\ 585; 631 f.). --- In these last places the
ground of spatial intuition is even quite predominantly ascribed to the
thing in itself.

From these intimations it is seen that Kant derived not only the material
but also the form from the thing in itself. But when this is so, it is
unjustified when he so often sets up a dilemma that offers the choice
between letting the object determine the representation and the
representation determine the object. In his own theory both are in reality
presupposed. Kant cannot therefore consistently be the strict apriorist he
will be. There is a fundamental presupposition that he has not brought
forward, and on which nevertheless his whole «system of principles» rests:
namely, that the thing in itself acts constantly. For if this is not the
case, then --- even if the nature of the subject be presupposed
unchangeable --- the constant elements in experience, which are the real
foundation for Kant's whole epistemological edifice, could change, and the
principles would then change too. --- From this side too it shows itself
that «the critical philosophy can only philosophise hypothetically». Kant's
assumption that he could carry through a philosophy about the phenomena
without needing to set up any presupposition whatever about the thing in
itself was incorrect. His first critics already drew attention to the fact
that when he uses the concepts thing, existence, and cause about it, it is
nevertheless not absolutely unknown. But it must be added that he also
presupposes that it acts in a regular manner.

Kant believed himself to have furnished a proof of the validity of the
causal proposition, something from which Wolff had come off most
unhappily. While Wolff would derive it from the principle of contradiction,
Kant grounds it as the presupposition of the possibility of
experience.\footnote{\textit{Kritik der reinen Vernunft.} 2nd ed.\ p.\
264.} Kant's proof rests on the assumption that the concept of cause is one
of the subjective forms without which experience in the strict sense does
not come into being. But just as little as the circumstance that we
apprehend outer phenomena in the form of space could be derived from the
nature of the subject alone, just as little does Kant regard it as
independent of the thing in itself that we apprehend phenomena as causes
and effects. This is seen from a statement such as this: «Dem
transscendentalen Objecte können wir allen Umfang und Zusammenhang unserer
möglichen Wahrnehmungen zuschreiben \ldots\ Die Erscheinungen sind, ihm
gemäss, \ldots\ gegeben».\footnote{Ibid.\ p.\ 522 f.} Here, then, the
phenomena are said, not only as regards material but also as regards form,
to be in agreement with the thing in itself, in that their connection too
is referred to it. The thing in itself must thus be just as constant in its
action on the subject as the connection among the phenomena is. And here we
stand, then, at a fundamental presupposition about the thing in itself
which Kant never formulated. It shows itself that to «the possibility of
experience», which bears Kant's whole theory of knowledge, there belongs
more than the subjective conditions. Hereby both apriorism and
phenomenalism are limited. The possibility of anticipating the phenomena
and the necessity of distinguishing between these and the thing in itself
are limited in a significant manner. In particular, the necessity falls
away of excluding the causal principle from validity for the thing in
itself. Although the strict, a priori proof for the causal proposition
falls away, there is gained the advantage that the great difficulty which
lay in Kant's assumption of the thing in itself as cause of the sensations
falls away at the same time. At the same time as the system is altered in an
empirical and realist direction, it also loses its great inconsistency.

\section*{III.\\ Theory and Practice}

\noindent 17.\ We return once more to Kant's standpoint in 1762 and the
following years, in order to consider it from a new side. Kant had seen the
necessity of a reformation of philosophy and expected salvation for the
present along the path of analysis. But when so much stood undecided and
unfinished, and when, with respect to some of the fundamental concepts
(e.g.\ the concept of cause), we stood before the inexplicable, then the
fund of theoretical philosophy that Kant at this time had at his disposal
for positive presentation cannot have been great. In his effort «rather to
have a thorough philosophy than a comprehensive one» he had made a clean
sweep. Even in the «Beweisgrund», the most constructive writing of this
time, he declares that the intention was not to give a real demonstration,
and that God's existence did not need to be proved either. In confidence
that religious faith had an entirely different foundation than
demonstration, he overturned the whole basis of natural theology. And in a
similar manner he proceeded in «Träume» against spiritualist psychology.
The cosmological antinomies seem at the same time to have played a great
role for him in his attempts to find the limits of knowledge. The material
for the part of the «Kritik der reinen Vernunft» that he called Dialectic,
and which was the part that made not the least impression at the work's
appearance, was thus in essentials finished.

Kant's philosophy was bound thereby to acquire a certain negative stamp, in
opposition to the dogmatism bursting with positivity. Kant was fully
conscious of this. He uses of his philosophy the expression «philosophy of
ignorance», also «negative philosophy». He finds this philosophy the most
difficult of all, because it must go back to the sources of knowledge in
order to be able to ground its non-knowledge.

«Where error is tempting and at the same time dangerous, there negative
cognitions and criteria are more important than positive ones and often
constitute the real object of our science \ldots\ Socrates had a negative
philosophy as regards speculation, namely concerning the worthlessness of
much supposed science and concerning the limits of our knowledge. The
negative part of education is the most important
(Rousseau).»\footnote{\textit{Reflexionen Kants.} II.\ p.\ 44 f.\ (Nos.\
148--151).}

These references to Socrates and to Rousseau are characteristic of Kant's
standpoint at this time. Socrates' knowledge was, after all, a knowledge
that he knew nothing, a knowledge grounded on clear insight into what was
required in order to know something. And Socrates also warned against
unfinished concepts and sought, by all-sided discussion, to lead to correct
determinations of concepts. As regards Rousseau, Kant draws in the cited
statement a parallel between his «negative education» and his own negative
philosophy. By negative education Rousseau understood «an education that
strives to perfect the organs, the instruments of our knowledge, before it
gives us knowledge itself, and which prepares for reason by exercising the
senses. Negative education does not fold its hands in its lap, far from it:
it imparts no virtues, but it prevents vices; it teaches not truth, but it
prevents errors.» It is grounded particularly by the necessity of coming to
know the nature and character of the individual, as they unfold of their
own accord, before one intervenes in accordance with general and handed-down
rules. It is, in its care to keep error out and to find the limits of the
individual's peculiar nature, a difficult art --- the art: \textit{tout
faire en ne faisant rien.}\footnote{\textit{Lettre à M.\ de Beaumont}
(Petits chefs d'œuvres de J.\ J.\ Rousseau.\ Paris 1859.\ p.\ 313). ---
\textit{Émile.} Livre II.\ (Paris 1851.\ p.\ 80 f.; 116.)} --- In several
respects the critical philosophy, according to the best and most justified
in its spirit, has its prototype in Rousseau's doctrine of negative
education. In both cases negation and critique serve to protect the
positive unfolding of life, whereas dogmatism, both in philosophy and in
pedagogy, is eager to intervene and to fix or regulate everything according
to its «eternal truths».

It is well known to what a high degree Rousseau's writings interested Kant.
When he got Émile, it kept him back from his usual walk. But had Kant's
notes and additions to his «Beobachtungen über das Schöne und Erhabene»
not at the last moment been rescued out of a grocer's pile of waste
paper,\footnote{\textit{Sämmtl.\ Werke.} Ed.\ Rosenkranz u.\ Schubert XI,
1.\ p.\ 218 f.} one would not have known how deeply personal this
expression was. It brought about for Kant an entirely new foundation for
the valuation of human beings and human relations. Hitherto Kant had been
an optimist, had regarded intellectual development as the highest and seen
progress decided by it. From Rousseau he learned a measure of human worth
that to a certain degree was independent of intellectual development. He
now learned that the rabble is not to be despised merely because «it knows
nothing». He «learned to honour human beings». And he praised Rousseau for
having brought forth man's deeply hidden nature beneath the forms of
civilisation, which all too often conceal it.\footnote{\textit{Sämmtl.\
Werke.} Ed.\ Rosenkranz u.\ Schubert XI, 1.\ p.\ 240.\ 248.} --- Here,
then, there was given a factual foundation for human life, whereby this
became independent of how things went with the scientific knowledge of
existence and its source. This was for Kant a great support, since he had
just grounded his critique of natural theology and of spiritualist
psychology. If the learned lost some favourite doctrines, then «the great,
to us most estimable multitude»\footnote{\textit{Kritik der reinen
Vernunft.} 2nd ed.\ Preface p.\ XXXIII.} lost nothing.

18.\ But it is with this Rousseauian influence as with the Humean one: were
it not secured by external testimony, we would in Kant's writings find no
compelling occasion to presuppose it. In itself one could very well, from
Kant's train of thought as it unfolded in the year 1762 and the following
years, have understood how he came to the distinction between theory and
practice, which from this time onward --- thus long before he fixed it in
his Critiques of Reason --- played so great a role in him.

With the critique of the proofs of God's existence the duality of the
motives that had led to the construction of the concept of God was bound to
come to light as well. The one group of motives was of a purely theoretical
kind and consisted essentially in the urge to bring the causal series to a
close, to let the working thoughts come to rest in the contemplation of a
being «grounded in itself». The other group was of an aesthetic-ethical
kind, and consisted essentially in an urge to concentrate in one concept
everything valuable in the world. By the fusion of both groups of motives
there arose the dogmatic philosophy's concept of God as \textit{ens
realissimum}, the sum-total of all reality, such that reality meant the
same as perfection, in that only such properties as express negation and
limitation were to be excluded. Already in the «Beweisgrund» (I, 4, 3) Kant
opposed this confusion of the concepts reality and perfection. Perfection
expresses a valuation, which is made by a feeling and willing being: it can
thus not be an absolute property. And in the treatise on negative
magnitudes (III, 2nd remark) Kant draws attention to the fact that, since
displeasure is just as positive and real as pleasure, one cannot make the
sum-total of all reality identical with the highest perfection. This from
the metaphysical side. From the psychological side the matter is treated in
the treatise on the distinctness of the principles, where it is said (IV,
2): «Only in our day has one begun to see that the capacity to discern the
true is cognition, the capacity to feel the good is feeling, and that they
must not be confused with one another. Just as there are unanalysable
concepts of the true, that is, of what is found in the objects of knowledge
considered in and for themselves, so too is there an unanalysable feeling
of the good, which is never found in a thing in and for itself, but always
in relation to a feeling being. Practical philosophy has its principles,
just as theoretical philosophy has its own, and their origin is still to be
investigated, toward which Hutcheson and others have made a good
beginning.» --- In «Nachricht von der Einrichtung seiner Vorlesungen» it is
stated that correct moral judgments can be passed by the human heart
without the circumlocution of demonstrations, by what one calls sentiment.
Reference is made to Shaftesbury, Hutcheson, and Hume as those who have come
furthest in seeking out the first grounds of morality. It is Rousseau's
predecessors that Kant here refers to. Of Rousseau himself we are reminded
both when he demands, as the foundation of ethics, a study of human nature,
which is easily mistaken on account of the changes that external relations
bring about, and when the task is set: to decide what perfection belongs to
man in the state of natural simplicity and in the state of wise
simplicity.\footnote{This opposition Kant treated far later, with special
reference to Rousseau's various writings, in his treatise: \textit{Muthmasslicher
Anfang des Menschengeschlechts} (1786).} But only because we know by other
means of Rousseau's influence on Kant do these traces become clear.

Kant was, as is seen, himself on a path that was bound to lead him to
sympathise with Rousseau. The latter's influence has probably, like Hume's,
been awakening, accelerating, has led to a quicker and more definite
unfolding of thoughts that were already budding, but has not consisted in a
positive imparting of something that did not lie as a possibility in his
own preceding development. Of not least importance has it been that
Rousseau taught that human beings, in the intellectual domain as well as in
the material, have many imagined needs. \textit{Nous pouvons être hommes
sans être savants!} --- \textit{Jeune homme, sachez être ignorant!} ---
says the Savoyard Vicar to his disciple. When Kant now stood at the break
with dogmatic speculation and had overturned the foundation for natural
theology and spiritualist psychology, he was bound to find a support in
Rousseau's train of thought. In the concluding section of «Träume eines
Geistersehers» this is clearly traced. After his critique of the
speculative and the visionary attempts to gain access to or insight into
the spirit-world he exclaims: «How much is there that I do not understand!»
--- but adds at once: «And how much is there that I do not need!»

The emancipation of ethics from metaphysics and theology was herewith
given. But Kant acknowledges not merely a foundation for the ethical
independent of theory. He goes a step further, in that he states that it
would even be very questionable if moral action were motivated merely by
the belief in another world. Instead of building morality on religion, one
should conversely build religion on morality. Those who would know all
about the other world must wait until they come there. For the present they
know that they are to do their duty, and more they do not need to know.
Kant ends his ingenious work as Voltaire ends his Candide: Let us go out
into the garden and work!

On this point too, then, there is continuity between the earlier works and
the later. The concluding section of «Träume» is developed further twenty
years afterward in «Kritik der praktischen Vernunft». And that the insight
gained was held fast in the meantime is seen from the fact that in the
Dissertation (§\ 9) a distinction is expressly drawn between the highest
theoretical and the highest practical concept.

19.\ With respect to the foundation on which Kant built the ethics
independent of metaphysics and theology, considerable variations took place
in the course of his development, though in such a way that there is a
certain fundamental stamp common to the various stages through which his
ethical conception passes.

In «Beobachtungen über das Schöne und Erhabene» [Observations on the
Feeling of the Beautiful and the Sublime] (1764 [i.e.\ 1763]) true virtue
is already made dependent on definite principles, in opposition to «adopted
virtue», which rests on immediate natural inclinations such as compassion
and benevolence. What is characteristic of this first ethical standpoint of
Kant's is that the principles on which true virtue rests are «the
consciousness of a feeling that lives in every human breast, \ldots\ the
feeling of the beauty and dignity of human
nature.»\footnote{\textit{Beobachtungen über das Gefühl des Schönen und
Erhabenen.} Königsberg 1764.\ p.\ 23.} By the principled emphasis on the
significance of the general principles Kant distinguishes himself from the
English moralists, on whom he otherwise relies in these years. But the
principles themselves are assumed to develop in that one becomes conscious
of one's feeling of the beauty and dignity of human nature. This feeling is
thus the moral feeling. In it, as is developed in «Träume eines
Geistersehers», the will of the individual is bound to the general will: «A
secret power compels us to direct our aim also toward the good of others or
according to a foreign will, even though it often happens unwillingly and
in high degree conflicts with the selfish inclination; --- and the point
where the directional lines of our drives converge thus lies not merely in
us, but there are also in the will of others outside us forces that move us
\ldots\ We see ourselves, then, in the most secret motives, dependent on
the rule of the general will, and from this there springs, in the world of
thinking beings, a moral unity and a systematic constitution according to
purely spiritual laws.»\footnote{\textit{Träume eines Geistersehers.} I,
2.\ (Kehrbach's edition, p.\ 23.)}

Kant's ethics here stands on purely psychological ground. What proclaims
itself in the individual mood and in the individual case is subordinated to
what is demanded by regard for the whole world to which the individual
belongs. The enlarged feeling, conditioned by the most circumspect
experience, we become conscious of in general principles and thereafter
control the emotion and urge of the moment. By this subordination of the
limited under the universal the ethical feeling acquires the character of
the sublime, and it is from this side that Kant comes to treat of it in
«Beobachtungen».

20.\ But Kant did not (unfortunately) remain at this connection of
psychology and ethics. His analysis led him, in the ethical domain as well
as in the epistemological, to distinguish sharply between reason and
sensibility and between form and material. It even seems as if this
distinction in the ethical domain was found by him earlier than in the
theoretical domain.\footnote{Letter to M.\ Hertz of 21 February 1772: «In
der Unterscheidung des Sinnlichen vom Intellektualen in der Moral und den
daraus entspringenden Grundsätzen hatte ich es schon vorher [i.e.\ before
the Dissertation of 1770] ziemlich weit gebracht.» Cf.\ letter to Herder, 9
May 1767, to Lambert 2 Sept.\ 1770.} The psychological-genetic conception,
which in «Beobachtungen» and «Träume» he still regarded as compatible with
ethics, he pushed aside, in order that the ethical might have its source in
pure reason. The influence of the dilemma he now sets up between reason and
sensibility is that all feeling is reckoned to the sensible part of our
nature,\footnote{«Alles Gefühl ist sinnlich.» \textit{Kritik der
praktischen Vernunft.} 1, 1, 3 (Kehrbach's edition, p.\ 92).} to the
material, the empirical, that by which we are passive, not self-active. In
diametrical opposition to the standpoint in «Beobachtungen», now the
principles are to determine the feeling, not conversely. The moral feeling
now becomes the feeling of respect determined by the consciousness of the
universal law. Respect for other personalities is grounded by the fact that
they are organs for the moral law, for the same law that we ourselves feel
within us. And while Kant in «Beobachtungen» named the feeling for the
beauty of human nature alongside the feeling for its dignity, this feeling
of beauty has later fallen away entirely from his ethics. This feeling
presupposed a harmony between the elements of human nature, which Kant, on
account of his sharp opposition between reason and sensibility, could no
longer assume. Therefore he was so doubtful about Schiller's juxtaposition
of «Anmuth und Würde» [grace and dignity]. Although it was his opinion that
duty should be exercised with joy and gladness, there must nevertheless be
no sensible elements involved therein. The gladness shall be conditioned
exclusively by the very unfolding of the inner power of the rational will,
not co-conditioned by the harmonious relation to the sensible nature. In
any case grace must not be placed before dignity; for even though the
consciousness of duty may have grace as its companion, it must nevertheless
not take any account of it. The mood that the consciousness of duty awakens
has the character of the sublime, not of the beautiful, and where ethical
work is concerned, the Graces must keep at a respectful distance. Only when
Hercules has subdued the monsters can he move among the Muses and the
Graces.\footnote{Kant's answer to Schiller's «Anmuth u.\ Würde» is found in
his «Religion innerhalb der Grenzen der blossen Vernunft». 2nd ed.\ p.\ 10
f. --- Cf.\ also notes written after the reading of Schiller's treatise in
«Thalia»: \textit{Lose Blätter} I.\ p.\ 120; 126 f.}

Sharp as the opposition is between Kant's earlier ethical standpoint and the
definitive one, we can nevertheless, with the help of the Kantian notes
published in recent years, trace the transition from the one standpoint to
the other.

21.\ When Kant, by his own statement, learned from Rousseau «to honour human
beings», and when he (in «Beobachtungen») conceived the moral feeling as
the feeling of humanity's dignity, on what, then, in his eyes, did this
honour and dignity rest? --- There emerges here, early in Kant's ethical
conception, a chief trait that never left it, although he explained it
philosophically in very different ways. It was the value he attached to the
capacity to subordinate the individual and changing under the general and
unchangeable, impressions and moods under the principle fixed in spite of
all vicissitudes, the passively received under spiritual activity. This
trait made him, already in «Beobachtungen», reckon the moral feeling under
the feeling of the sublime, just as he also later (see both «Kritik der
praktischen Vernunft» and «Kritik der Urtheilskraft») found a close kinship
between these two feelings, though in such a way that he no longer reckoned
the moral feeling under the feeling of the sublime, but rather took the
reverse path, in that he interpreted the feeling of the sublime by assuming
a co-operation of the moral feeling.

For the understanding of the change in Kant's ethical theory, which seems to
have taken place in the interval between 1766 (Träume) and 1770
(Dissertation), perhaps the following consideration may serve. --- The more
clearly the relation of opposition between the individual, changing, passive
on the one side, and the universal, fixed and active on the other side, was
emphasised, the more there had to become, for Kant, something unsatisfactory
in conceiving the principles as having arisen through our becoming conscious
of the moral feeling. They thereby seem to become dependent on
contingencies. A quite certain starting point for the ethical would by
contrast be obtained if the activity of reason were posited as the first and
original, and feeling were regarded only as its product. Thus the relation
between principle and feeling had to be reversed if the direction in which
Kant's ethics already moved in «Beobachtungen» and «Träume» was to be
carried through absolutely. And that this tendency had to be strengthened by
Kant's untiring work at distinguishing between «the pure» and «the
empirical» follows of itself. It was perhaps not even free of the fact that
the concept of the pure acquired a moral tinge in Kant's eyes. At least
there is an ambiguity in the manner in which he from this time onward uses
the expression «Sinnlichkeit» of all elements in human nature with the
exception of reason.

He now distanced himself from the English moral philosophers, to whom he had
earlier referred, and regarded it as a great fault that they traced the
moral criteria back to the feeling of pleasure and displeasure. He saw no
difference in principle between Shaftesbury's and Epicurus' standpoint
(Dissertatio §\ 9), something over which Mendelssohn (letter to Kant 23
Dec.\ 1770) already rightly complained. The moral concepts are now to be
known not by experience, but by pure reason, and moral philosophy belongs,
as regards the first principles of judging (\textit{principia dijudicandi
prima}), to pure philosophy (Diss.\ §\ 7; 9). --- How Kant at this stage
more closely conceived the grounding of the ethical is not known. Yet there
are statements and drafts that show that prior to the ethical theory he
developed in «Grundlegung zur Metaphysik der Sitten» and in «Kritik der
praktischen Vernunft» there goes a theory that to a certain degree forms a
transitional link between the older psychologising and the later
rationalist ethics.

22.\ In the «Reflexionen» there is found, from a moment that does not lie far
from the genesis of the «Dissertation», even if it lies, in the editor's
view, prior to it, a determination of the concept of practical idealism,
according to which this consists in the fact that happiness does not depend
on the external world, or more fully: «that the world is only what we make
it, that in a glad mind it has a bright, in a melancholy mind a dark
appearance, that the grounds of a happy condition are to be sought in
ourselves.» And this practical idealism is expressly called «not the
idealism of fancy, but of reason, of wisdom.»\footnote{\textit{Reflexionen
Kants.} II.\ p.\ 319.\ (Nos.\ 1116; 1119.)}

Not a little akin to the practical idealism thus determined is the
conception that comes forth in a moral-philosophical draft in the «loose
leaves of Kant's papers» published by R.\ Reicke. Here morality is defined
as «the idea of freedom as the principle of happiness». «It is true that
virtue has the advantage that it would bring about the greatest welfare by
means of what nature offers. But its high value consists not in the fact
that it, as it were, serves as a means. That it is we ourselves who produce
it as authors without regard to the empirical conditions (which can give
only particular maxims of living), --- that it brings self-contentment with
it, --- this is its inner value». «Only he is in a position to be happy
whose use of his will does not conflict with the conditions for happiness
that nature gives him \ldots\ Happiness is really not the greatest sum of
pleasure, but pleasure in the consciousness of my own power to be content
(\textit{Selbstmacht zufrieden zu sein}); it is at least the most important
formal condition for happiness, though there are still others required
\textit{materialiter} (as by experience)». «Happiness must derive from an a
priori ground that reason approves». «Freedom is a capacity to do and to
refrain independently of empirical grounds \ldots\ But I am free only from
the compulsion of sensibility, not from the limiting laws of reason \ldots\
The unboundedness with which I can will what is against my very will, so
that I cannot count on myself with certainty, must in the highest degree
awaken displeasure in me. It is known a priori that a law is needed by which
freedom is limited in such a way that the will can agree with itself. This
law I cannot betray without conflicting with my reason, which alone can fix
the practical unity of the will according to principles». «It is freedom
according to the laws of thoroughgoing agreement with itself that
constitutes the value and dignity of personality.»\footnote{\textit{Lose
Blätter} I.\ p.\ 10--15. This fragment was by the editor, the librarian
Rudolph Reicke, assigned to the years 80 or 90. Since it was clear to me
from the content that it must date from a time prior to «Grundlegung zur
Metaphysik der Sitten», which appeared in 1785, I applied to the editor for
information concerning the grounds that had determined his dating of the
manuscript. Mr.\ Reicke has been so good as to investigate the matter anew
and has, in a private letter of 10 Nov.\ 1892, informed me that there is
nothing in the way of the fragment's dating from the 1770s, but that he
would nevertheless, by the handwriting, rather assign it to the 1780s; to
the 1790s he now considers it impossible to assign it. Possibly it then came
into being at the end of the 1770s or the beginning of the 1780s. I would be
inclined to believe that it was composed some time before the last redaction
of «Kritik der reinen Vernunft», since it does not yet have the pronounced
imperative ethics that «Kritik d.\ r.\ V.» already hints at, any more than
the doctrine of the intelligible character.}

The ethical standpoint that lies at the foundation of these statements
differs from that which appears in «Beobachtungen» by its rationalist and
individualist character. It finds the highest in the activity by which the
individual can be the creator of his own happiness, and finds as the
essential law for this activity that the individual must act with full
consistency, in agreement with himself. There is here made an essential
step in the direction of laying the chief weight on the formal side of the
ethical. The ethical law shall abstract from the content, spring from the
free activity itself, and can then be only its
self-limitation\footnote{The fragment's expression about the limitation of
freedom might recall «Kritik der reinen Vernunft» 1st ed.\ p.\ 569, where
there is talk of «the principle by which reason sets limits to the freedom
that is in itself lawless», and where the Stoic sage is mentioned as an
ethical ideal.} through the necessity of preserving the agreement with
itself. Kant had here reached a chief point in his ethics: the inner
relation between action and law, a relation which, in his conception, could
become the right one only when the law was purely formal; only then could it
be a priori, abstract from all definite content. In genetic respect,
therefore, that moral-philosophical fragment is exceedingly interesting. A
problem it forms by its eudaemonistic and individualist character, which
places it in opposition both to the earlier and to the later standpoint. The
relation between virtue and happiness, which for Kant's later ethics became
so difficult a problem that only religious postulates could resolve it,
contained, according to this «practical idealism», no difficulty in
principle. The external goods must be given up, or one must in any case
prepare oneself for resignation toward them; but this is outweighed by the
feeling of freedom and power that here stands as the highest good, just as it
is determined by what here is the highest virtue. The standpoint in the
fragment forms a swinging-arc in Kant's moral-philosophical development of an
analogous kind to «the Dissertation» in his epistemological development. Of
greatest importance in both places was the sharp distinction between form and
content, a priori and a posteriori. It became determining for the future.
The sharper form in which this opposition appears in Kant's definitive
ethics brought with it, among other things, also that opposition between
virtue and happiness which could be resolved only along the path of the
religious postulate.

If one would imagine what, from the standpoint of the fragment, may have led
Kant over to his definitive ethics, then it lies near at hand to refer to
how he, in his theoretical philosophy, had reached the point of determining
the universally valid as that which agrees with the nature of all rational
subjects. By analogy herewith the moral law became not merely an expression
of the individual's agreement with himself, but of the preservation of the
connection of reason in the whole, of the individual's connection with the
whole intelligible world. Thus it is said in a note from the 1770s: «The
intelligible world has laws according to which I am fit for any world, not
merely for this one, for the world of sense, but no matter how my nature and
external nature, as well as the society in which I live, are arranged \ldots\
It accords with the rules of prudence when this [i.e.\ prudence] is
extended.»\footnote{\textit{Reflexionen Kants.} II.\ p.\ 336.\ (No.\
1158.)} These last words seem to hint at a development from the earlier
rationalist standpoint to the definitive one. By means of the idea of the
intelligible world Kant has, despite the formalism, succeeded in giving his
rationalist ethics a social character such as his first,
psychologically-motivated ethics had by its very genesis.

But yet another, quite different train of thought has probably also
co-operated in this transition, and that a train of thought which precisely
had a psychological-historical character and so far might seem to accord
little with the rationalist path on which Kant had entered. In the beginning
of the 1780s Kant was strongly occupied with the thought that the human
natural endowments, insofar as they concern the use of reason, come to full
development only in the species, not in the individual, since the development
of reason requires far longer time than the individual's life-span. The goal
for human development can therefore not be an individual one, but must be a
universal goal, a goal for all, not merely a goal for this or that
individual. This train of thought lies at the foundation of Kant's
remarkable treatise «Idee zu einer allgemeinen Geschichte in
weltbürgerlicher Absicht», which appeared in 1784, the year before
«Grundlegung zur Metaphysik der Sitten». In the cited train of thought one
must in part seek the explanation of the fact that Kant now conceives the
moral law not as the formal rule within the individual's isolated striving,
but as the rule that brings the individual striving into agreement with the
striving of all rational beings. It is, unconsciously for Kant himself, an
idea drawn from psychological-historical experience that has given his
definitive ethics its universalist character alongside the rationalist
character. The very universal formula that Kant makes into the moral
principle contains, as it were, a recollection of a society of individuals
whose wills it is a matter of bringing into harmony with one another; it is
really an idealisation of the laws of intercourse among personal beings that
experience shows us.\footnote{Cf.\ what I have already remarked about this in
my treatise «Om Grundlaget for den humane Etik» (Copenhagen 1876).\ p.\ 52.}
The great value of Kant's formula rests precisely --- however much Kant
himself would deny it --- on this its crypto-empirical origin.

23.\ Finally, renewed investigation and analysis of the immediate ethical
feeling has also exercised its great influence in the genesis of Kant's
definitive ethics. Just as analysis and construction could be shown to
co-operate in the grounding of Kant's theory of knowledge, so we also find
both methods in the grounding of his ethics. It is precisely the
after-effect of this immediate contact with the ethical in practice that
gives Kant's ethical writings (especially «Grundlegung zur Metaphysik der
Sitten») so fresh and vigorous a colour. In «Grundlegung» Kant precisely
makes his way to the proper moral philosophy through analysis of «the common
moral knowledge of reason» --- little as he otherwise will let psychological
experience have any influence whatever on the ethical ideas. By this renewed
analysis he then found, as a chief trait of the ethical, that the individual
feels himself bound in his inner being to a law that is at once the
expression of his own innermost essence and orders him as a member of a
great realm of reason. The moral criterion now becomes whether the maxim
that the action expresses can be made into a universal maxim. This doctrine
forms an analogy to the doctrine of the law-governed connection of phenomena
as the expression of real experience.

To the law corresponds the force. From the moral law as a fact proclaiming
itself in all consciousness of duty Kant infers the pure self-activity of
reason, which he calls freedom (in one of the several meanings in which he
unfortunately uses this word). For freedom, self-activity, cannot be
immediately apprehended, and just as little can it be derived from
experience. The moral law we become immediately conscious of, and it leads
to the assumption of freedom as the capacity to act independently of all
sensible conditions.\footnote{\textit{Kritik der praktischen Vernunft,} §\ 6
remark (Kehrbach's ed.\ p.\ 53). Cf.\ \textit{Lose Blätter} I, p.\ 120. ---
When in individual places, e.g.\ Grundlegung p.\ 108, it is said that «that
in man which is pure activity comes immediately to consciousness», the word
«immediately» must, according to the context, be understood as opposed to
«durch Afficirung der Sinne».} The moral law is the expression of reason's
autonomy.\footnote{«Das moralische Gesetz ist das Gesetz der Autonomie der
reinen practischen Vernunft». \textit{Kritik der praktischen Vernunft.}
Kehrbach's ed.\ p.\ 53.} By obeying the law of his own essence the individual
at the same time obeys the general law of the spiritual world.

Herewith Kant's ethics differs from the earlier rationalist ethics, as it
was developed by the «Intellectualists» of the 17th century and by Price.
When these grounded the ethical on reason, they meant thereby the objective,
rational connection of things, which man apprehends by his reason and
thereafter bows before. Reason as a human faculty here had, as Fr.\ Jodl
aptly puts it,\footnote{\textit{Geschichte der Ethik in der neueren
Philosophie.} I.\ Munich 1882.\ p.\ 221.} a merely receptive role to carry
out, namely to appropriate the reason that was assumed to be given in the
relations of the world, and to apply it in one's action. Kant's practical
reason is by contrast, like his theoretical reason, an expression of man's
own innermost essence, a sum-total of laws for man's spiritual activity. The
profound thought in Kant's definitive ethics is that man, by becoming
conscious of the moral law, becomes conscious of himself as a citizen of two
different worlds that meet in his person: «The moral ought is the subject's
own willing as a member of an intelligible world and is thought only as an
ought insofar as the subject at the same time regards itself as a member of
the world of sense.»\footnote{\textit{Grundlegung zur Metaphysik der
Sitten.} 3rd ed.\ Riga 1792.\ p.\ 113.} In strict rationalist form Kant has
here developed anew a thought that came forth in more empirical form in
«Beobachtungen», when appeal was made to the dignity of human nature. It is
this dignity for which he believed he had to seek a deeper grounding than
psychology could give. The idea of man's dignity is the continuous element
in Kant's ethics; what varies is the grounding.

Despite its exalted character, Kant's rationalist ethics nevertheless rests,
when considered as a member of his critical philosophy, on a great
inconsistency. The parallel between the theoretical and the practical reason
is not maintained. The enemy whom Kant combats with the weapons of critique
he sees as highly different in the two domains. In the theoretical domain it
is speculation, the transcendent use of the concepts, that is to be combated.
«Kritik der reinen Vernunft» corresponds to its name, for it is really a
critique of reason. But «Kritik der praktischen Vernunft» is not at all a
critique of reason's application in the ethical-practical domain, as one
would suppose from the title and from the analogy with that principal work.
And Kant really admits it himself. The task here is, he says (Einleitung von
der Idee einer Critik der praktischen Vernunft), «to keep the empirically
conditioned reason away from the presumption of wishing to provide
exclusively the determining ground of the will.»\footnote{Cf.\ more closely
on this: A.\ Fouillée: \textit{Critique des systèmes de Morale
contemporains.} Paris 1883.\ p.\ 129--142.} By contrast he does not discuss
the presumption of pure reason to wish to be the exclusive determining
ground, but assumes precisely that the ethical consists solely in this. He
does not admit that there could be a critical investigation to undertake
concerning reason's justification in the practical domain, and that
notwithstanding that he himself strongly emphasises the inexplicability in
the fact that pure reason determines the will of a «sensibly» conditioned
being. --- It avenges itself here that the psychological foundation has been
pushed aside.

\section*{IV.\\ The Copernican Principle}

\noindent 24.\ I return to Kant's theoretical philosophy in order to
investigate how the turning point in 1769, which brought about the proper
grounding of the critical philosophy, can be inserted as a member in Kant's
continuous development, and how the consequences were gradually drawn from
this turning point.

What had been reached in 1762 was the recognition that the concepts with
which one had hitherto operated in philosophy were useless, unfinished;
analysis was not carried through; of philosophical construction there could
only, in a distant future, be any talk. Kant had not given up the hope of a
concluding theoretical philosophy that gave information about questions
lying beyond the domain of experience. But he turned with mockery and irony
against the dogmatists, who, unperturbed by all difficulties, built their
systems and decided all problems. His attention was directed toward method
and toward the determination of the limits of knowledge, and he had already
begun to speak of a critique of reason. In order to find the limits of
rational knowledge he operated with the method of antinomies, in that in
the contradiction between two trains of thought each grounded for itself he
saw a proof that reason had ventured too far out. «I sought», he
says,\footnote{\textit{Reflexionen Kants.} II.\ p.\ 4 (No.\ 3).}
«something that was certain, if not with respect to the objects, then at
least with respect to the nature and limits of knowledge. I found
gradually that many of the propositions we regard as objective are in
reality subjective, that is, contain the conditions under which alone we
discern or comprehend the objects.» From the content of consciousness Kant
thus came to turn his attention to conscious activity itself and to make it
the object of his analysis. Of the procedure in this analysis there has
been talk above (§\ 6).

It is through such an analysis that he came not only to distinguish between
form and material, but also between intuiting and thinking. Thus it is said
in «Prolegomena» (p.\ 119): «In an investigation of the pure (containing
nothing empirical) elements of human knowledge it succeeded for me only
after long reflection to distinguish and single out the pure elementary
concepts of sensibility (space and time) from the elementary concepts of
the understanding.» --- At the same time there dawned on him the conviction
that in those elementary concepts of sensibility we have only an expression
for the forms in which, according to the nature of our sense-faculty, we
apprehend things, not for the things' own essence, and that it is by this
alone that mathematical knowledge of nature is possible. These were the
important steps taken in 1769 and developed in the Dissertation of 1770.

25.\ From the year before there is a little interesting treatise by Kant:
«Vom ersten Grunde des Unterschiedes der Gegenden im Raume» [On the First
Ground of the Distinction of Regions in Space]. He here defends Newton's
conception, according to which absolute space lies at the foundation of
every individual determination of place, against Leibniz's conception of
space as a property or a relation determined by the nature and action of
things. He relies especially on the fact that with bodies entirely alike
and yet incongruent (such as, for example, the right and left hand) one
cannot determine the place of the parts without taking account of the space
outside the bodies. From this he infers that space as a whole lies at the
foundation in the determination of the individual place.

In the course of the next year (1768--69) it must then have dawned on Kant
that «absolute space» was not anything objective and real, but a subjective
form, a schema for our apprehension of things. Here, then, we have an
example of the transposition from objective reality to subjective condition
of which he speaks in the utterance cited above. It is in this connection of
interest to note that the same argument which Kant in the little treatise of
1768 used to prove Newton's conception of space against Leibniz's, he later
uses (Prolegomena §\ 13)\footnote{See also \textit{Metaphysische
Anfangsgründe der Naturwissenschaft.} Riga 1786.\ p.\ 7 f.} to prove his
own doctrine of the subjectivity of space against the usual assumption of
its objectivity. It is then clear that Newton's conception had been for him
a station on his path. It did not lie so far off, for one who, to so high a
degree as Kant, had occupied himself with the study of Newton, to make the
transposition from the objective to the subjective. Newton conceived space
as \textit{sensorium dei}, a thought that is seen to have occupied Kant
much.\footnote{See on this Benno Erdmann in «Reflexionen Kants». II.\ p.\
104 f.\ (the note).} It was then merely a matter of putting \textit{sensorium
hominis} in the place of \textit{sensorium dei}, which, for one occupied
with investigating the general nature and mode of operation of human
knowledge, could likewise not lie so far off. --- The after-effects of
Newton's conception one finds in the psychological difficulties in which
Kant's theory of space entangles itself by retaining the determination of
space as infinite and yet conceiving it as the form of the sense-intuition
operative in experience, whose object must always be finite ---
difficulties that Kant himself comes to underscore by speaking of space as
a given infinity.\footnote{See on this B.\ Erdmann in «Reflexionen Kants».
II.\ p.\ 110 f.\ (the note). --- Kant himself has, in a treatise directed
against Kästner, which was first brought forward by Dilthey in \textit{Archiv
für Gesch.\ der Phil.} III (p.\ 83 ff.), sought to clarify this point.}

The demonstration by which Kant seeks to establish space and time as
subjective forms of intuition consists partly in direct analysis, partly in
an apagogic proof. For the sake of brevity I speak in the following only of
space, since Kant's reasoning for time runs quite parallel to his reasoning
about space.

Space is a form of intuition, because it presents to us a whole that
becomes conscious to us by successive grasping-together of the parts, by
synthesis. A concept of the understanding, by contrast, shows us merely the
composition of the whole as conditioned by the parts, without showing us its
coming-into-being. For intuition the individual part is conditioned by the
whole, for the understanding conversely (Diss.\ §\ 1; 15 Coroll.). --- If,
on the other hand, space were not a form of intuition but a concept of the
understanding, then the concepts of continuity and infinity would enclose a
contradiction, which falls away when one supports them on intuition's
constant and in principle unceasing advance from part to part (Diss.\ §\ 1;
2, III).

A double demonstration --- a direct and an apagogic --- is also employed to
establish that space is only a subjective mode of apprehension.

If we apprehend space as a subjective schema in which we order the
sensations that arise in us (\textit{schema coordinandi}), then the laws of
space must also hold for all material phenomena. For everything we are to be
able to apprehend must become an object for the ordering mental activity
(\textit{vis animi, omnes sensationes coordinans}). The laws of
sense-intuition must therefore necessarily be the laws of nature, and the
application of geometry in natural science is thereby explained: the same
intuition from which geometry springs we use, after all, every time we
apprehend something by means of the outer sense. The validity of applied
geometry is then a simple inference from the subjectivity of the form of
space (Diss.\ §\ 15 C.\ and E.). --- On the other hand: «If the concept of
space were not originally given by means of the nature of consciousness,
then the application of geometry in natural science would be uncertain; for
doubt might then be raised as to how far this very concept, drawn from
experience, agreed with nature» (Diss.\ §\ 15 E.).

Kant's demonstration is of interest, since it shows us how direct analysis
and antinomic trains of inference go hand in hand. After he has found by
analysis that space is a synthetic form of intuition, he shows that another
assumption would lead to absurdities. And as he then passes over to drawing
the further inferences from what he has thus established, he investigates the
two possibilities: that space is subjective, and that it is not subjective,
and shows that the consequences of both lead to the same result, namely that
the possibility of applied mathematics is bound to the subjectivity of the
forms of intuition. We find here the conclusion of the long series of
analyses and antinomies through which Kant has worked his way forward to the
decisive principle for his theory of knowledge.

This principle we may, employing the well-known comparison in the preface to
the 2nd edition of the «Kritik der reinen Vernunft», call the Copernican
principle. It states that our knowledge of existence is determined by a
nature of knowledge, and that we, when we know its laws, can decide
beforehand which general laws the phenomena must be subject to. --- This
principle did not --- which cannot surprise --- dawn on Kant at once in its
full range. In the Dissertation it is applied only to sense-intuition, not
to the understanding. While sense-intuition thus shows us things only as
phenomena, that is, as they appear to our apprehension, which clothes them
in the forms of our grasping intuition, the understanding is still, in its
pure concepts (substance, cause, possibility, actuality, necessity, etc.),
to be able to lead to a knowledge of things in themselves, of the inner
essence of existence. It had, by Kant's own statement, cost him long
reflection to establish the difference between intuition and understanding,
which dogmatic philosophy had effaced. It was then no wonder that he could
believe that these two different branches of our faculty of knowledge were
subject to different conditions, so that subjectivity and the restriction to
phenomena held only for intuition, not for the understanding. The venerable
distinction, deriving from Plato, between \textit{Phænomena} and
\textit{Noumena}, to which he had now been led along his own
path,\footnote{The expression «intelligible world» is already found in
«Träume» (1.\ Th.\ 2.\ Hauptst.), where it is used of the «mystical
philosophy's» world of spiritual substances, evidently with a view to
Leibniz's world of monads. Leibniz himself used the expression \textit{mundus
intelligibilis} of the world of monads (Epistola ad Hanschium. Op.\ phil.\
ed.\ Erdmann, p.\ 445 C.), or, what comes to the same, of the world of
purposes (Animadv.\ in Principia Cartesiana. Phil.\ Schr.\ ed.\ Gerhardt.\
IV, p.\ 391). --- B.\ Erdmann and R.\ Riedel have shown, what is fully
confirmed by the «Reflexionen», that Kant's intelligible world is Leibniz's
world of monads with due modification.} was bound, for him, quite naturally
to correspond to the difference between intuition and understanding. There
is, however, nothing new in the knowledge of noumena that he still regards
as possible along the path of thinking. It consists essentially in the
thought, familiar from «Beweisgrund» and earlier works, of existence's
unitary ground as the necessary condition for the law-governed interaction
in nature. (See above §\ 4.) The Dissertation thereby hangs together
continuously with Kant's earlier train of thought, while at the same time,
by its Copernican principle, it points forward toward «Kritik der reinen
Vernunft».

Friederich Paulsen, who places the awakening in 1769, conceives the
Dissertation as a reaction, brought about by Hume's influence, against the
sceptical tendencies of the preceding years. I believe that he sets it in
too great opposition to the immediately preceding period. Kant had never
given up the hope of a concluding result along the path of thinking, though
its fulfilment was deferred to an indeterminate future. He had even, in his
most sceptical writing, declared himself in love with metaphysics, and in a
contemporary letter stated that it was merely a matter of divesting
metaphysics of its dogmatic garb. It is psychologically understandable that
he, after the discovery of the difference between intuition and understanding
and of the Copernican principle in its application to sense-knowledge, could
believe himself near the goal. For a moment he believed he had reached the
definitive determination of limits that he had sought after for several
years. Add to this that the knowledge of noumena which Kant at this stage
regards as possible is yet only symbolic, since all data of intuition are
lacking. However strictly Kant distinguishes between intuition and
understanding, he nevertheless impresses already now (Diss.\ §\ 10) that all
material for our knowledge is obtained by sensation. And the concepts of the
understanding are conceived no less than the forms of intuition as an
expression of the nature of the knowing mental activity (Ibid.\ §\ 8). This
is rather a limitation of the noumenal knowledge than an extension of it.

26.\ What Kant lacked in order to reach his definitive standpoint was the
carrying-through of the Copernican principle on all the first presuppositions
for our knowledge, and thereby the extension of phenomenalism to hold for all
scientific knowledge of existence. It might seem as if there were not far to
go, since the first step, the one that usually costs the most, had already
been taken with the grounding of the critical doctrine of space and time. But
it nevertheless cost Kant eleven years' work to take the step all the way. A
sign of the man's great conscientiousness, and at the same time a testimony
to how highly different the matter stands for the thinker stuck in the midst
of the problems and for the retrospective historian. It can already on
practical grounds be understood that Kant shrank from definitively excluding
knowledge from the thing in itself. The carrying-through of phenomenalism
became possible for him only when his ethical theory had developed in such a
way that in the ethical basic fact he could see a possibility of access to the
world that was now closed off for theoretical knowledge. The «unconditioned»
that he could no longer accommodate in knowledge he had to find a place for
in faith; only then did he regard his task as solved, in that he had found
«practical data» for a concept he could not give up, but whose theoretical
data he had dissolved.\footnote{See on this \textit{Kritik der reinen
Vernunft.} 2nd ed.\ p.\ XIX--XXII, where these «practical data» are regarded
as a verification of «the Copernican principle», just as Newton's law of
gravitation gave a verification of Copernicus' astronomy.} The chief
difficulty has, however, as is seen from the «Reflexionen», been theoretical.
It was, after all, a matter of finding a point of view for the activity of the
understanding by which what had been shown for sense-intuition could also come
to hold for it, despite the difference so strongly emphasised in the
Dissertation.

Yet the proper Kantian philosophy came into being with the Dissertation. So
Kant himself conceived it, when he (in the letters to Garve and Mendelssohn)
declared «Kritik der reinen Vernunft» to be the fruit of twelve years' work,
and when he set the year 1770 as the limit for the writings he wished
reprinted. (See above §\ 1.) I cannot see that the assumption maintained by
Benno Erdmann, that the awakening first took place in the middle of the
1770s, is compatible with this. It was, after all, a matter in these years of
carrying through an already found principle, a principle that in itself
presupposes a break with dogmatism. Kant cannot, when in 1783 he wrote the
well-known lines in the preface to the Prolegomena that were to cause Kant
scholars so much head-scratching, have wished to divide into two sharply
separated parts the period that he otherwise regarded as a whole. The
definitive closing of access to the things in themselves is, to be sure, an
epoch-making event. But the expression awakening does not fit it, whereas it
does fit a change in the manner of inquiry, which brought about that closing
as its final result.

27.\ In the Dissertation a distinction is drawn between a merely logical use
of the understanding, which consists in comparison, the demonstration of
identity and difference, and a real use, which is not more closely
characterised, but which is to have absolute validity, not merely validity
for the phenomena, whereas the logical use is merely to serve the ordering of
the phenomena. The work of thought that lies between the Dissertation and
«Kritik der reinen Vernunft» aims at a closer investigation of how matters
stand with this «real use».\footnote{The expression «real use of reason», in
opposition to the merely logical use, also occurs in «Kritik der reinen
Vernunft», at least in the Dialectic. (2nd ed.\ p.\ 355.)} The problem posed
itself (as is seen from Kant's letters to Lambert and Hertz) immediately
after the publication of the Dissertation. When the concepts of the
understanding or «intellectual representations» that we use to penetrate, along
the path of thinking, into the essence of things, express our mind's mode of
operation, how can they then be shown to have validity for existence itself?
Or as it is said in the letter to Hertz of 21 February 1772: «When such
intellectual representations [Diss.\ §\ 8 names: the concepts possibility,
actuality, necessity, substance, cause] rest on our inner activity, whence
then comes the agreement that they are to have with objects, which after all
are not produced by them, and how do the principles of pure reason about the
objects agree with these, when this agreement may not seek support in
experience?»

As guidance in this investigation Kant had the discovery from the
Dissertation. There can, as regards the mathematical fundamental relations,
be given an a priori knowledge of sensible nature, because the mathematical
fundamental concepts express the laws for the synthetic construction that
displays itself in every sense-intuition. Now it is clear that so long as Kant
held fast his earlier so strongly maintained conception of thinking as
analysis, the real use of the understanding had to stand as a riddle. It had
in its time advanced Kant a great step that he so definitely impressed that
all judging was a comparison. (The treatise on false subtlety.) Now it was a
matter of gaining a deeper-reaching conception of the activity of the
understanding, so that this could offer an agreement with sense-intuition. To
regard the understanding straightforwardly as an intuition would lead to
mysticism. What way out, then, was to be found?

In the «Reflexionen» and in the «loose leaves» one can see how diligently
Kant worked with the various fundamental concepts that constituted the content
of the earlier ontology, before he got this so-called science transformed into
«the analytic of the understanding». It was first a matter of tracing them
back to certain main concepts, in order then in these to find expression for
definite kinds of activity of the understanding. Already in the cited letter
to Hertz he mentions his effort «to trace all the concepts of absolutely pure
reason back to a certain number of categories, in such a way that their
division into classes is determined by a few fundamental laws holding for the
understanding». Here there were two paths to go. One could proceed from an
analysis of the fundamental concepts that the understanding operates with, in
order to find the laws for the understanding's activity, or one could seek
directly to determine the law for the understanding's activity, in order then
thereby in turn to find which main kinds of fundamental concepts it applies,
according to its nature, to the given material. Kant has, as it seems, gone
both ways.

First he tried his way with various combinations and divisions of the
«ontological» concepts. He has, for example, dwelt on the concepts
comparison, connection, and relation (or connection), and he has particularly
investigated the concepts connection and comparison and found that by both
the operations they express a unity is brought about or established: «All
unity is either the unity of comparison or of connection. The first we have
insofar as something is of the same kind as something else; the second
insofar as a manifold is connected in one ground.»\footnote{\textit{Reflexionen
Kants.} II.\ p.\ 143 (No.\ 461). Cf.\ p.\ 164 (Nos.\ 525 ff.).} He has
investigated the concepts substance, cause, and interaction, which evidently
particularly interested him. And here too it was the concept of a unity in
the manifold, brought about by the form of thought, that he particularly
noted: «The concept of substance and accident itself contains a synthesis,
likewise cause and effect, and manifoldness in real unity [i.e.\
interaction]. According to its various relations to the inner sense, nature
must now necessarily stand under one of these syntheses.»\footnote{\textit{Lose
Blätter.} I.\ p.\ 18 f.\ Cf.\ \textit{Reflexionen Kants.} II.\ p.\ 174--181
(Nos.\ 562--587).} The concept synthesis he seems first to have applied only
to these three relations, which he early gathered under the expression
relation. Thus it is said in a note that precedes the genesis of the proper
doctrine of categories: «Only for the relation do objectively synthetic
propositions about the phenomena hold.»\footnote{\textit{Lose Blätter.} I.\
p.\ 17. --- The expression relation Kant has probably taken from the
Aristotelian table of categories. He applied it to the above-mentioned group
of categories before he used it to designate a class of judgments. See on
this Adickes: \textit{Kants Systematik als systembildender Faktor.} Berlin
1887.\ p.\ 26; 36. --- Perhaps he was, by the concept of relation, led to fix
his attention on the judgments. Cf.\ Refl.\ No.\ 532: «Our reason contains
nothing but relations». No.\ 596: «The category of relation
(consciousness-unity) is the most important of all. For unity really
concerns only relation. Thus this constitutes the content of the judgments in
general and lets itself alone be thought as determined a priori».} It cannot
(according to what has been developed in the first chapter of this treatise)
surprise us that it was particularly this class of categories that led Kant
to extend the concept synthesis to the activity of the understanding and
thereby also to discover that the Copernican principle had also to hold for
this. --- Already in the beginning of the Dissertation he had intimated, as a
common point of view for all our knowledge, that it partly forms a whole out
of given elements, partly dissolves given wholes into their elements. Thereby
arise two limit-concepts: the absolutely simple (\textit{simplex} i.e.\
\textit{pars quæ non est totum}) and the world-totality (\textit{mundus}
i.e.\ \textit{totum quod non est pars}). Here, then, reference was made to
the two fundamental mental operations, synthesis and analysis, and the
question was bound to present itself, which of them was the more primitive. In
the Dissertation it was now established, as regards intuition, that it rested
on synthesis. The great advance that Kant made in the psychology of knowledge
in the course of the 1770s was then the discovery that not only all
intuiting, but also all thinking consists in synthesis, when one goes back to
the fundamental form of the activity of thought, and that analysis is always
secondary in relation to it, since one can only dissolve what is already
connected. It dawned on Kant that not only the ordering in space and time,
but also the inner connection among the phenomena which our knowledge brings
about, is due to a «linking-together» (\textit{Verkettung}) of
representations, an inner act of consciousness which first brings a real
whole out of the given elements. Without such an act of consciousness no
experience; thus it must have universal validity.\footnote{\textit{Lose
Blätter.} I.\ p.\ 16.}

The Copernican principle stood, so far as one can see, firm for Kant in its
full extent as soon as he had, by investigation of the categories, found
synthesis as the common determination of all cognitive activity. Our mind,
our self, then came to stand as an archetype or prototype for all objects
that we are to know. (See above §\ 5 concl.; 12; 15.)

But investigation of the current ontological fundamental concepts and
demonstration of their origin from a synthetic mental activity was not the
only path by which Kant came to his doctrine of categories. It was for him
only a provisional method, which did not satisfy him, because it was too
empirical. It contained no guarantee that all the fundamental concepts had
been found. And for Kant it was a chief matter that the limits of knowledge
could be marked out a priori, after a complete investigation of all the
fundamental forms of knowledge. This he believed he could attain by laying
the doctrine of judgments at the foundation, in such a way that to every
particular kind of judgment there corresponds a certain definite category.
But before he struck out upon this path, which as is known led him to encumber
his presentation with a great scholastic scaffolding, the general idea of the
activity of the understanding as synthesis had clearly enough stood clear for
him. This is seen from what has been cited above from the «loose leaves», and
the «Reflexionen» also establish it. Thus when it is said (Refl.\ No.\ 600):
«The judgment is the unity of consciousness of the manifold in the
representation of an object in general. The category is the representation of
an object in general, insofar as it is determined with respect to this
objective unity of consciousness». Here the concept synthesis lies at the
foundation for the definition of the judgment. It must have been found as a
fundamental concept before it could be seen that all judging falls under it.
The presentation in «Kritik der reinen Vernunft» testifies to the same. Here
the understanding is first determined as the faculty of knowing by means of
concepts. A knowledge by concepts is discursive, thus presupposes a mental
function, and a function means «the unity of the act of ordering various
representations under a common representation» [cf.\ the Dissertation's
\textit{vis coordinandi}]. Of his concepts the understanding can now make use
only by judging, that is, determining the representation of an object by
another representation; it is only intuition that has immediately to do with
the object.\footnote{\textit{Kritik der reinen Vernunft.} 1st ed.\ p.\ 67--69
(Von dem logischen Verstandesgebrauche überhaupt).} The functions of the
understanding can be completely found when one can completely present the
various ways in which unity of various representations is expressed in the
judgments. And this Kant now believes himself able to accomplish by an
«improved» edition of the current logical doctrine of
judgments.\footnote{On the manner in which Kant, for his «transcendental»
use, has altered the system of judgments, see Adickes: \textit{Kants
Systematik,} p.\ 30--41.} On this I shall not enter more closely. As
significant as Kant's general thought about the activity of the understanding
as judging and about the judgment as synthesis is, just as arbitrary is the
application in detail, although his great genius shows itself here too in many
valuable thoughts to which he finds room within the scholastic frame. Had
Kant, instead of laying so much weight on deducing a complete system of
categories and principles, held fast to the general thought that the logical
principles that condition the formal validity of the judgments must also hold
for all experience, for all connection among phenomena that is to come into
being for us, then he would have avoided much obscurity and arbitrariness,
and the transcendental analytic, as the doctrine of the possibility of
applied logic, would then have been placed in an instructive manner alongside
the «transcendental aesthetic» as the doctrine of the possibility of applied
mathematics.

28.\ The «subjective deduction» (that is, the derivation, supported on
psychological analysis, of the nature of the faculty of knowledge), by which
Kant arrived at determining synthesis as the fundamental form of knowledge and
at giving a survey of the special forms in which it displays itself, was for
him nevertheless not the chief matter. His chief task was by contrast the
objective or transcendental deduction, the demonstration of the validity and
limits of rational knowledge. To this the analytic investigation of the forms
forms only introduction and foundation. It was, as he wrote to Garve (7 August
1783), his task to ground «an entirely new and hitherto unattempted science,
namely the critique of an a priori judging reason», which «from the mere
concept of a faculty of knowledge (when it is exactly determined) derives a
priori all objects, everything one can know about them, indeed even what one
involuntarily, though wrongly, is compelled to judge about them». With
justice he complained that he was so grossly misunderstood on this point in
the first assessment of his principal work, and this had as a consequence
that the psychological analysis stepped more into the background in his later
presentations. (See above §\ 13.) The attractive thing about the first
edition of «Kritik der reinen Vernunft» consists precisely in the freshness
and fullness with which the concept of synthesis is presented in its
significance for the cognitive functions, all the way from sense-intuition to
the highest act of thought. It is a question whether this presentation, which
in the 2nd edition is considerably abbreviated, is not the part of Kant's
doctrine that is most penetrating and has the most enduring value. As so
often, it is not the conscious chief aim, but the subordinate aims, which one
had to set oneself in order to reach the former, that have proved to be of
greatest importance.

The significance of Kant's concept of synthesis is, in the first place, that
he thereby shows us what it is to understand a thing. His doctrine of
knowledge rests on the simple thought that to understand something is to see
it in as close a connection as possible with everything else we know. A
connecting mental activity must make itself felt in all knowledge; only when
we can apply it have we appropriated the given material and feel the peculiar
satisfaction that accompanies the relieving of a need. What we do not
understand is the isolated and the connectionless. Naturally there are many
degrees between the two extremes: to understand and not to understand,
connection and connectionlessness, and, as we have seen (above §\ 14--15),
Kant committed the error of not attending to these differences of degree. But
the merit he has is to have been the first to see the problem of knowledge in
this simple way. It is a Columbus' egg that he has found and stood on its end.
Remarkably enough Hume had found the same Columbus' egg, but let it roll away
from him. In his «Treatise» (which Kant did not know) Hume really came to the
same fundamental concept that Kant reached in «Kritik der reinen Vernunft».
For Hume there traced the problem of knowledge from the special problems
(particularly the problem of the real validity of geometry and of the causal
relation) back to the fundamental difficulty which for him --- on account of
the atomistic psychology from which he proceeded --- lay in the connection
among representations as a whole. The great riddle lay for him at last in the
principle of unity (\textit{the uniting principle}).\footnote{\textit{Treatise
of human nature.} I, 3, 14 (ed.\ Selby-Bigge.\ Oxford 1888.\ p.\ 169). Cf.\
Appendix (ibid.\ p.\ 636).} Kant now shows that the difference between the
intelligible and the enigmatic rests on the difference between what can be
grasped together in continuous connection and that where such a function is
not possible. The connecting function itself one can then plainly only find
enigmatic, either when one proceeds from the assumption that the
connectionless should be the intelligible, or when one asks about the greater
connection by which our grasping function is in turn conditioned. The first
path leads out into the meaningless; the other path leads to posing for us
the last problem that the theory of knowledge and psychology as a whole can
stand before, and which marks their limit: namely, of the arising of
knowledge itself, or of consciousness as a whole, in existence. Thought is to
itself its own last problem.\footnote{Cf.\ my \textit{Psychology.} 3rd ed.\
p.\ 403.} --- It would surely have been of great importance for Kant's
development if he had known Hume's principal work at an early stage of his
career. The collision between the two great thinkers would then have
concerned a still more central question, and Kant would perhaps have been led
to a complete working-through of the psychological foundation for his theory
of knowledge, while on the other hand he would have been able to spare himself
a long scholastic labour. The egg would earlier and more easily have been
stood on its end. It was, for the subsequent history of philosophy, a fateful
act that Hume committed when he, dissatisfied with the slight literary fortune
that his youthful work had and vexed at the inopportune attacks to which it was
subject from the theological side, disowned this magnificent work of
his.\footnote{Eduard Grimm (\textit{Zur Geschichte des Erkenntnissproblems.}
Leipzig 1890.\ p.\ 584) sees in this disowning an admission on Hume's part of
the one-sidedness in the «Treatise». But that the motive was decidedly the
one brought forward above is clearly seen partly from Hume's autobiography,
partly from the letter in which Hume urged his publisher to let the
declaration in which he disowned the «Treatise» accompany the «Inquiry».
\textit{Letters of David Hume to William Strahan.} Edited by C.\ Birkbeck
Hill.\ Oxford 1888.\ p.\ 288 f. --- It was also for the last-mentioned reason
that Hume did not let the «Dialogues» appear in his lifetime. Ibid.\ p.\
330.}

While Kant by his diligent work of thought found the fundamental concept that
could lead beyond the atomistic psychology that lay at the foundation of
empiricism, he had by the same fundamental concept attained a conception of
the essence of mind that led beyond the surreptitious concepts of the
spiritualist dogmatism (cf.\ above §\ 8). Over against empiricism, which would
regard the unity of mind as a mere result of the manifold impressions, he
maintains the unitary activity as the fundamental stamp of spiritual life,
which cannot be explained by external influence alone; over against
spiritualism, which has indeed an eye for this fundamental stamp, but
dogmatically will trace it back to a mystical substance behind consciousness,
he maintains (in his critique of rational psychology) that from the function
one has no right to infer the substance, and that this inference in any case
cannot lead to real knowledge.\footnote{Leibniz is the one of the dogmatic
philosophers who comes closest to the concept of synthesis, as when he
(Monadologie §\ 14) defines «perception» thus: \textit{l'état passager qui
enveloppe et représente une multitude dans l'unité ou dans la substance
simple.} --- Leibniz reproached the Cartesians for not defining what they
understood by thinking (i.e.\ consciousness). His own definition was that
consciousness is manifoldness expressed in a unity. (\textit{Philos.\ Schr.}
ed.\ Gerhardt.\ VII.\ p.\ 551). «\textit{Les âmes sont des unités et les
corps sont des multitudes, mais les unités ne laissent de representer les
multitudes \ldots\ C'est dans cette réunion que consiste la nature admirable
du sentiment.}» (ibid.\ p.\ 542.)} Kant's concept of synthesis expresses the
fundamental form and fundamental law for spiritual life, so far as we can know
it along the path of experience. He has thereby given psychology an important
point of view, a measure in the valuation of the development of spiritual
life, --- and at the same time some of its most important problems, for the
spiritual unity is often threatened by the manifoldness and conflict of the
elements. It is especially interesting to see how high Kant, by this concept,
raises himself above the Enlightenment's emphasis on the clearly conscious and
intellectually seen-through. For synthesis lies indeed at the foundation of
consciousness, but need not be noticed by consciousness itself. It is
consciousness's constant condition, but not necessarily its object. It works
blindly, instinctively in our innermost essence. When Kant in «Kritik der
Urtheilskraft» (§\ 46) defines genius as «the innate mental endowment by which
nature gives rules to art», then it is only a special application of something
that holds for all conscious life. There is, if one will, something of genius
in every conscious being, which in those who are especially to be called
geniuses only steps forth in the highest potency. --- The immediate and
involuntary in spiritual life is thus, by the Kantian conception, maintained
in its right and significance as the constant foundation. Kant's concept of
synthesis has a far more comprehensive significance than the one he ascribes
to it in his theory of knowledge, and it has in any case more enduring worth
than other components of his philosophy which he himself rated ever so highly.

\end{document}
